% c11-wasserstoff/de.tex — Kapitel 11, Das Wasserstoffatom (deutsche Leitfassung)
%
% Smoke B der Phase 0.  Mathematik aus evidence/derivations/c11/DERIVATION.md
% (Herleitung) und REVIEW.md (unabhängige Prüfung); jede Aussage hat ihre Zeile
% in evidence/claims.json (c11.cl.*). scripts/c11/check_c11_formulas.py prüft exakt
% die in formulas.json gelisteten Zeilen ungesternter equation/align-Umgebungen,
% deren linke Seite eine Ylm-, Rnl-, psinlm- oder Lag-Funktion mit numerischen
% Indizes und Argumentliste ist (Lag darf einen symbolischen oberen Index haben).
% Gesternten Displays, symbolischen Definitionen, Integralen, CDFs und tpcdef-Inhalten
% gilt diese Prüfung nicht. Gelistete Gleichungen und formulas.json werden zusammen geändert.

\chapter{Das Wasserstoffatom}\label{c11}

Das Wasserstoffatom besteht aus einem Elektron und einem Proton. Es ist das einfachste
Atom. In diesem Kapitel konstruieren wir eine Familie gebundener Zustände in einem
quantenmechanischen Modell: Ihre Energien und Wellenfunktionen sowie die unter einer
zusätzlichen Photonenergie-Konvention verglichenen Wellenlängen stehen als geschlossene
Formeln da.

Wir gehen dabei in drei Schritten vor. Zuerst schreiben wir das Atom als das hin, was es
ist, nämlich als System aus \emph{zwei} Teilchen. Dann spalten wir die Bewegung des
Schwerpunkts ab; diese Rechnung führen wir vollständig aus, denn sie entscheidet, was die
berühmten Orbitalbilder eigentlich zeigen. Übrig bleibt ein Problem für ein einziges
Teilchen mit der sogenannten reduzierten Masse, und dieses lösen wir.

\section{Zwei Teilchen, ein Atom}\label{c11.s1}

Das Elektron sitzt am Ort $\posE$, das Proton am Ort $\posP$; beide Orte sind Vektoren
im dreidimensionalen Raum. Die beiden Teilchen haben die Massen $\massE$ und $\massP$
und ziehen einander elektrisch an. Mehr steckt in unserem Modell nicht:

\tpcdef{D60}

Für Wasserstoff ist $\Znuc = 1$. Wir lassen $\Znuc$ in den Formeln stehen; sie gelten
dann ebenso für einen Kern mit der Ladung $\Znuc\echarge$ und einem einzigen Elektron,
wenn man für $\massP$ die Masse dieses Kerns einsetzt.

\begin{tpcfigure}{c11.fig.zweikoerper}
  \caption{Die Orte $\posE$ des Elektrons und $\posP$ des Protons, vom Ursprung $O$ aus
    gemessen, und die beiden neuen Koordinaten aus Abschnitt~\ref{c11.s2}: die
    Schwerpunktskoordinate $\posCM$ und die Relativkoordinate $\posRel = \posE - \posP$.
    Der Schwerpunkt teilt die Verbindungsstrecke im Verhältnis der Massen. Das Bild ist
    für das Massenverhältnis $\massP : \massE = 3 : 1$ gezeichnet, damit man den
    Schwerpunkt vom Proton unterscheiden kann; das ist nicht das Verhältnis im
    Wasserstoffatom.}
\end{tpcfigure}

% GAP-TJO: Die Formulierung der Brücke vom endlichdimensionalen Kern (C0-C4) zu L^2 ist
% TJOs (PLAN.md Abschnitt 14, Punkt 7; sie gehört nach C5).  Der folgende Absatz hält sich
% an die Aufzählung dort und erfindet nichts hinzu.
Was heißt hier „Zustand“ und was „Operator“? Ab Kapitel~\ref{c05} arbeiten wir mit
Wellenfunktionen, und in genau diesem Sinn ist das Folgende gemeint: Die Zustände, mit
denen wir rechnen, sind durch eine Wellenfunktion $\Psi(\posE, \posP)$ gegeben, eine
komplexwertige Funktion von sechs reellen Variablen, deren Betragsquadrat sich über den
ganzen $\mathbb{R}^6$ zu eins integriert; das Skalarprodukt zweier Wellenfunktionen ist
ein Integral; und wir verwenden zunächst den Differentialausdruck aus \defref{D60}
auf Funktionen, die sich außerhalb der Kollisionspunkte zweimal differenzieren lassen.
Zweimalige Differenzierbarkeit allein legt keinen Operator-Definitionsbereich auf
$\Ltwo{\mathbb{R}^6}$ fest. Die Eigenfunktionsgleichungen gelten außerhalb der
Kollisionspunkte und als $L^2$-Identitäten fast überall. Auf welchen Funktionen genau
die Vorschrift einen Operator festlegt, ist eine eigene Frage\zoomref{c11.f1}.

Zu zwei Teilchen gehören also Funktionen von sechs Variablen. Man schreibt das als
\begin{equation}\label{c11.eq.hilbertraum}\claim{c11.cl.hilbertraum}{source-quoted}
  \Ltwo{\mathbb{R}^3} \otimes \Ltwo{\mathbb{R}^3} = \Ltwo{\mathbb{R}^6},
  \qquad
  (f \otimes g)(\posE, \posP) = f(\posE)\, g(\posP).
\end{equation}
Links steht das Tensorprodukt, mit dem man in der Quantenmechanik Systeme zusammensetzt;
rechts steht, was das für Wellenfunktionen bedeutet: Das Produkt $f \otimes g$ einer
Funktion des Elektronenorts und einer Funktion des Protonenorts ist die gewöhnliche
Produktfunktion \srcref{reed-simon-1}{Theorem II.10 (a), S. 52}. Für uns ist nur die
rechte Seite wichtig: Wir rechnen mit Funktionen von sechs Variablen.

Unser Modell lässt vieles weg, und das sagen wir gleich dazu. Elektron und Proton sind
Punktladungen ohne Spin; die Bewegung ist nichtrelativistisch; die einzige
Wechselwirkung ist die Coulombanziehung; ein Lichtfeld, das das Atom aussenden oder
aufnehmen könnte, kommt in $\op{\HamTot}$ nicht vor. Alle Ergebnisse dieses Kapitels sind
Aussagen über dieses Modell.

\section{Schwerpunkt und Relativbewegung}\label{c11.s2}

Das Potential in $\op{\HamTot}$ hängt nur vom Abstandsvektor $\posE - \posP$ ab. Das legt
nahe, diesen Vektor als eine der neuen Koordinaten zu nehmen. Als zweite nehmen wir den
Schwerpunkt.

\tpcdef{D61}

\tpcdef{D62}

\begin{tpcfigure}{c11.fig.massen}
  \caption{Wo der Schwerpunkt auf der Strecke vom Proton zum Elektron liegt: nach \defref{D61} im Abstand $(\massE/\massTot)\,\rad$ vom Proton. Bei gleichen Massen ist das die Mitte, beim Massenverhältnis aus Abbildung~\ref{c11.fig.zweikoerper} ein Viertel der Strecke. Im Wasserstoffatom ist das Proton so viel schwerer als das Elektron, dass der Schwerpunkt im Maßstab dieses Bildes mit dem Proton zusammenfällt.}
\end{tpcfigure}

\subsection{Die neuen Koordinaten}\label{c11.s2.1}

Zuerst lösen wir nach den alten Koordinaten auf. Aus $\posRel = \posE - \posP$ folgt
$\posE = \posP + \posRel$. Eingesetzt in
$\massTot\posCM = \massE\posE + \massP\posP$ ergibt das
$\massTot\posCM = \massTot\posP + \massE\posRel$, also
\begin{equation}\label{c11.eq.umkehr}\claim{c11.cl.umkehr}{derived}
  \posP = \posCM - \frac{\massE}{\massTot}\,\posRel,
  \qquad
  \posE = \posCM + \frac{\massP}{\massTot}\,\posRel.
\end{equation}
Zu jedem Paar $(\posE, \posP)$ gehört also genau ein Paar $(\posCM, \posRel)$ und
umgekehrt (Abbildung~\ref{c11.fig.zweikoerper}). Eine Wellenfunktion können wir deshalb
ebenso gut als Funktion der neuen Koordinaten schreiben. Wir nennen diese Funktion
$\widetilde{\Psi}$:
\begin{equation}\label{c11.eq.tilde}\claim{c11.cl.tilde}{derived}
  \widetilde{\Psi}(\posCM, \posRel)
  = \Psi\Bigl(\posCM + \frac{\massP}{\massTot}\,\posRel,\;
              \posCM - \frac{\massE}{\massTot}\,\posRel\Bigr),
  \qquad
  \widetilde{\Psi} = \op{\Uni}\Psi.
\end{equation}
Der Buchstabe $\op{\Uni}$ steht für die Abbildung, die aus $\Psi$ die Funktion
$\widetilde{\Psi}$ macht. Sie ändert nichts an der Physik, nur an der Beschriftung.

Ändert sie die Normierung? Bei einem Wechsel der Integrationsvariablen kommt es auf den
Betrag der Determinante der Koordinatentransformation an. Die Transformation mischt die
drei Raumrichtungen nicht: In jeder Richtung gilt dasselbe, zum Beispiel für die
$x$-Komponenten $x_{\mathrm{e}}$, $x_{\mathrm{p}}$ der alten und $X$, $x$ der neuen
Koordinaten
\begin{equation}\label{c11.eq.jacobi}\claim{c11.cl.jacobi}{derived}
  \begin{aligned}
    \begin{pmatrix} X \\ x \end{pmatrix}
    &= \begin{pmatrix} \massE/\massTot & \massP/\massTot \\ 1 & -1 \end{pmatrix}
      \begin{pmatrix} x_{\mathrm{e}} \\ x_{\mathrm{p}} \end{pmatrix}, \\
    \det \begin{pmatrix} \massE/\massTot & \massP/\massTot \\ 1 & -1 \end{pmatrix}
    &= -\frac{\massE + \massP}{\massTot} = -1.
  \end{aligned}
\end{equation}
Für alle drei Richtungen zusammen ist die Determinante $(-1)^3 = -1$. Ihr Betrag ist
eins: Der volle Jacobi-Faktor erhält das Integrationsvolumen und damit die Normierung.
Die besondere ebene Masche in Abbildung~\ref{c11.fig.koordinaten} hat denselben
Flächeninhalt wie das dort gezeigte Einheitsquadrat.
\begin{equation}\label{c11.eq.norm}\claim{c11.cl.norm}{derived}
  \int \abs{\widetilde{\Psi}(\posCM, \posRel)}^2 \,\dd^3\posCM\,\dd^3\posRel
  = \int \abs{\Psi(\posE, \posP)}^2 \,\dd^3\posE\,\dd^3\posP .
\end{equation}

\begin{tpcfigure}{c11.fig.koordinaten}
  \caption{Der Koordinatenwechsel für eine Raumrichtung, wieder für
    $\massP : \massE = 3 : 1$ gezeichnet. Durchgezogene Linien: $X$ fest; gestrichelte
    Linien: $x$ fest. Die neuen Koordinatenlinien stehen nicht senkrecht aufeinander,
    aber die Masche zwischen $X = 1$, $X = 2$, $x = 0$ und $x = 1$ hat denselben
    Flächeninhalt wie das Einheitsquadrat der alten Koordinaten: Der Betrag der
    Determinante in \eqref{c11.eq.jacobi} ist eins.}
\end{tpcfigure}

\subsection{Die kinetische Energie in den neuen Koordinaten}\label{c11.s2.2}

Jetzt kommt die eigentliche Rechnung. Wir brauchen die zweiten Ableitungen nach den
alten Koordinaten, ausgedrückt durch Ableitungen nach den neuen. Das Werkzeug ist die
Kettenregel, und wieder genügt eine Raumrichtung. Nach \defref{D61} ist
$X = (\massE x_{\mathrm{e}} + \massP x_{\mathrm{p}})/\massTot$ und
$x = x_{\mathrm{e}} - x_{\mathrm{p}}$, also
\begin{equation}\label{c11.eq.partiell}\claim{c11.cl.partiell}{derived}
  \frac{\partial X}{\partial x_{\mathrm{e}}} = \frac{\massE}{\massTot},
  \qquad
  \frac{\partial x}{\partial x_{\mathrm{e}}} = 1,
  \qquad
  \frac{\partial X}{\partial x_{\mathrm{p}}} = \frac{\massP}{\massTot},
  \qquad
  \frac{\partial x}{\partial x_{\mathrm{p}}} = -1.
\end{equation}
Eine Funktion der neuen Koordinaten hängt über $X$ und über $x$ von $x_{\mathrm{e}}$ ab.
Die Kettenregel sammelt beide Abhängigkeiten ein:
\begin{equation}\label{c11.eq.kette}\claim{c11.cl.kette}{derived}
  \frac{\partial}{\partial x_{\mathrm{e}}}
  = \frac{\massE}{\massTot}\frac{\partial}{\partial X} + \frac{\partial}{\partial x},
  \qquad
  \frac{\partial}{\partial x_{\mathrm{p}}}
  = \frac{\massP}{\massTot}\frac{\partial}{\partial X} - \frac{\partial}{\partial x}.
\end{equation}

\begin{tpcfigure}{c11.fig.kette}
  \caption{Die Kettenregel \eqref{c11.eq.kette} als Schema. Eine Funktion von $X$ und $x$ hängt auf zwei Wegen von $x_{\mathrm{e}}$ ab (durchgezogen) und auf zwei Wegen von $x_{\mathrm{p}}$ (gestrichelt). An jedem Pfeil steht der Faktor aus \eqref{c11.eq.partiell}; die Ableitung nach einer alten Koordinate ist die Summe über ihre beiden Wege.}
\end{tpcfigure}

Das wenden wir zweimal an. Wenn die Funktion zweimal stetig differenzierbar ist, kommt
es auf die Reihenfolge der Ableitungen nicht an, und man darf rechnen wie bei einer
binomischen Formel:
\begin{equation}\label{c11.eq.zweite}\claim{c11.cl.zweite}{derived}
  \begin{aligned}
    \frac{\partial^2}{\partial x_{\mathrm{e}}^2}
    &= \frac{\massE^2}{\massTot^2}\frac{\partial^2}{\partial X^2}
     + \frac{2\massE}{\massTot}\frac{\partial^2}{\partial X\,\partial x}
     + \frac{\partial^2}{\partial x^2}, \\
    \frac{\partial^2}{\partial x_{\mathrm{p}}^2}
    &= \frac{\massP^2}{\massTot^2}\frac{\partial^2}{\partial X^2}
     - \frac{2\massP}{\massTot}\frac{\partial^2}{\partial X\,\partial x}
     + \frac{\partial^2}{\partial x^2}.
  \end{aligned}
\end{equation}
Für die $y$- und die $z$-Richtung gilt dasselbe. Die Summe über die drei Richtungen
macht aus den zweiten Ableitungen Laplace-Operatoren; die gemischten Terme summieren
sich zu $\nabla_{\posCM} \cdot \nabla_{\posRel}$, dem Skalarprodukt der beiden
Gradienten:
\begin{equation}\label{c11.eq.laplace}\claim{c11.cl.laplace}{derived}
  \begin{aligned}
    \Lap_{\mathrm{e}}
    &= \frac{\massE^2}{\massTot^2}\,\Lap_{\posCM}
     + \frac{2\massE}{\massTot}\,\nabla_{\posCM} \cdot \nabla_{\posRel}
     + \Lap_{\posRel}, \\
    \Lap_{\mathrm{p}}
    &= \frac{\massP^2}{\massTot^2}\,\Lap_{\posCM}
     - \frac{2\massP}{\massTot}\,\nabla_{\posCM} \cdot \nabla_{\posRel}
     + \Lap_{\posRel}.
  \end{aligned}
\end{equation}
In $\op{\HamTot}$ steht $\Lap_{\mathrm{e}}$ geteilt durch $\massE$ und $\Lap_{\mathrm{p}}$
geteilt durch $\massP$. Genau diese Gewichte lassen die gemischten Terme verschwinden:
\begin{equation}\label{c11.eq.koeff}\claim{c11.cl.koeff}{derived}
  \begin{aligned}
    \frac{1}{\massE}\,\Lap_{\mathrm{e}} + \frac{1}{\massP}\,\Lap_{\mathrm{p}}
    &= \frac{\massE + \massP}{\massTot^2}\,\Lap_{\posCM}
     + \Bigl(\frac{2}{\massTot} - \frac{2}{\massTot}\Bigr)\,
       \nabla_{\posCM} \cdot \nabla_{\posRel}
     + \Bigl(\frac{1}{\massE} + \frac{1}{\massP}\Bigr)\,\Lap_{\posRel} \\
    &= \frac{1}{\massTot}\,\Lap_{\posCM} + \frac{1}{\massRed}\,\Lap_{\posRel}.
  \end{aligned}
\end{equation}
Vor $\Lap_{\posCM}$ steht der Kehrwert der Gesamtmasse, vor $\Lap_{\posRel}$ der Kehrwert
der reduzierten Masse aus \defref{D62}. Multipliziert man mit $-\hbar^2/2$ und nimmt das
Potential dazu, das nur von $\posRel$ abhängt, so steht das Ergebnis da.

\begin{theorem}[Abspaltung der Schwerpunktsbewegung]\label{c11.thm.separation}\claim[proof-included]{c11.cl.separation}{derived}
  Die Funktion $\Psi$ der beiden Orte sei zweimal stetig differenzierbar, wo immer
  $\posE \neq \posP$ ist, und $\widetilde{\Psi} = \op{\Uni}\Psi$ sei dieselbe Funktion in den
  neuen Koordinaten wie in \eqref{c11.eq.tilde}. Dann gilt an allen diesen Stellen
  \begin{equation}\label{c11.eq.separation}
    \begin{aligned}
      \op{\Uni}\,\op{\HamTot}\Psi &= \op{\HamCM}\widetilde{\Psi} + \op{\HamRel}\widetilde{\Psi}, \\
      \op{\HamCM} &= -\frac{\hbar^2}{2\massTot}\,\Lap_{\posCM},
      \qquad
      \op{\HamRel} = -\frac{\hbar^2}{2\massRed}\,\Lap_{\posRel} + \Vpot(\posRel).
    \end{aligned}
  \end{equation}
  Dabei wirkt $\op{\HamCM}$ nur auf die Schwerpunktskoordinate und $\op{\HamRel}$ nur auf die
  Relativkoordinate.
\end{theorem}
\begin{proof}
  Die kinetische Energie in $\op{\HamTot}$ ist $-\hbar^2/2$ mal der linken Seite von
  \eqref{c11.eq.koeff}; die rechte Seite von \eqref{c11.eq.koeff} liefert die beiden
  Laplace-Operatoren in \eqref{c11.eq.separation}. Das Potential
  $\Vpot(\posE - \posP)$ ist in den neuen Koordinaten $\Vpot(\posRel)$. Die Punkte mit
  $\posE = \posP$, also $\posRel = 0$, sind ausgenommen, weil das Potential dort nicht
  erklärt ist.
\end{proof}

\begin{tpcfigure}{c11.fig.zerlegung}
  \caption{Theorem~\ref{c11.thm.separation} als Schema; die Längen im Bild bedeuten nichts. Links das Atom, wie es dasteht: zwei Teilchen, gekoppelt durch das Potential. Rechts dasselbe System in den neuen Koordinaten: ein freies Teilchen der Masse $\massTot$ und, unabhängig davon, ein Teilchen der Masse $\massRed$ im Potential $\Vpot$.}
\end{tpcfigure}

Der Satz ist eine Aussage über Differentialausdrücke. Mit dem Tensorprodukt aus
\eqref{c11.eq.hilbertraum}, jetzt für die Variablen $\posCM$ und $\posRel$, schreibt man
dasselbe kurz als $\op{\Uni}\,\op{\HamTot}\,\op{\Uni}^{-1} = \op{\HamCM} \otimes \Id + \Id \otimes \op{\HamRel}$
\srcref{reed-simon-3}{§\,XI.4, S. 55}.

Der Hamiltonoperator zerfällt also in zwei Teile, die nichts voneinander wissen. Das
nutzen wir mit einem Produktansatz aus. Ist $\widetilde{\Psi}(\posCM, \posRel) =
\Phi(\posCM)\,\psi(\posRel)$, so wirkt $\op{\HamCM}$ nur auf $\Phi$ und $\op{\HamRel}$ nur auf
$\psi$:
\begin{equation}\label{c11.eq.produkt}\claim{c11.cl.produkt}{derived}
  (\op{\HamCM} + \op{\HamRel})\,\Phi\,\psi
  = (\op{\HamCM}\Phi)\,\psi + \Phi\,(\op{\HamRel}\psi).
\end{equation}
Der erste Summand gehört zu einem freien Teilchen der Masse $\massTot$. Der Schwerpunkt
des Atoms bewegt sich also, als gäbe es das Innenleben nicht. Der zweite Summand gehört
zu \emph{einem} Teilchen der Masse $\massRed$ im Potential $\Vpot$. In $\op{\HamRel}$ steckt
die ganze Atomphysik dieses Kapitels.

\begin{merksatz}[Exakt, nicht genähert]
  Die Zurückführung des Zweiteilchenproblems auf ein Einteilchenproblem mit der
  reduzierten Masse $\massRed$ ist in unserem Modell exakt. Eine Näherung wird daraus
  erst, wenn man $\massRed$ durch $\massE$ ersetzt, also so tut, als wäre das Proton
  unendlich schwer.
\end{merksatz}

Wie gut wäre diese Näherung? Das Elektron ist viel leichter als das Proton. Das
Verhältnis der Massen ist $\massE/\massP = \num[c11.cl.q]{5{,}44617 \cdot 10^{-4}}{}$.
Für die reduzierte Masse folgt daraus $\massRed/\massE = 1/(1 + \massE/\massP)$, in
Zahlen \num[c11.cl.mu-ratio]{0{,}999456}{}. Sie ist also nur wenig kleiner als die
Elektronenmasse (Abbildung~\ref{c11.fig.reduziert}). Weil die Rechnung mit $\massRed$ aber um nichts schwerer ist
als die mit $\massE$, rechnen wir durchweg mit $\massRed$ und sagen bei jeder Größe
dazu, welche Masse in ihr steckt.

\begin{tpcfigure}{c11.fig.reduziert}
  \caption{Die reduzierte Masse nach \defref{D62}, in Einheiten der Elektronenmasse, als Funktion des Massenverhältnisses; die waagerechte Achse ist logarithmisch. Bei gleichen Massen ist $\massRed$ die Hälfte von $\massE$. Je schwerer das zweite Teilchen, desto näher kommt $\massRed$ der Elektronenmasse; der rechte Punkt ist Wasserstoff.}
\end{tpcfigure}

\section{Was die Relativwellenfunktion ist und was nicht}\label{c11.s3}

Die Funktionen $\psi(\posRel)$, die wir in den nächsten Abschnitten ausrechnen, sind
Funktionen der Relativkoordinate. Bevor wir das tun, klären wir, was eine solche
Funktion über das Atom im Labor sagt. Dazu brauchen wir die zweite Hälfte des Zustands,
die Wellenfunktion $\Phi(\posCM)$ des Schwerpunkts.

Wir betrachten einen Produktzustand
$\widetilde{\Psi}(\posCM, \posRel) = \Phi(\posCM)\,\psi(\posRel)$, in dem $\Phi$ und
$\psi$ jede für sich normiert sind. Nach \eqref{c11.eq.norm} ist dann auch $\Psi$
normiert. Integriert man das Betragsquadrat über eine der beiden neuen Koordinaten,
bleibt die Dichte der anderen übrig:
\begin{equation}\label{c11.eq.relativdichte}\claim{c11.cl.relativdichte}{derived}
  \int \abs{\widetilde{\Psi}(\posCM, \posRel)}^2 \,\dd^3\posCM = \abs{\psi(\posRel)}^2,
  \qquad
  \int \abs{\widetilde{\Psi}(\posCM, \posRel)}^2 \,\dd^3\posRel = \abs{\Phi(\posCM)}^2.
\end{equation}
% GAP-TJO: Wie der Kurs operational sagt, was |psi|^2 "ist" (Wahrscheinlichkeitsdichte
% wofür, gemessen wie), ist TJOs Formulierung; hier steht nur die Randdichte der Formel.
Also ist $\abs{\psi(\posRel)}^2$ die Wahrscheinlichkeitsdichte für den Abstandsvektor
$\posRel = \posE - \posP$ vom Proton zum Elektron, und $\abs{\Phi(\posCM)}^2$ die für
den Ort des Schwerpunkts. Keine der beiden ist die Dichte des Elektrons im Labor. Diese
definieren wir eigens.

\tpcdef{D70}

\begin{tpcfigure}{c11.fig.kopien}
  \caption{Wie die Dichte des Elektrons im Labor zustande kommt, als Schema; Längen und Strichstärken bedeuten nichts. Jeder mögliche Ort des Schwerpunkts (Kreuze) trägt eine skalierte Kopie der Relativdichte (Ringe), mit $g_b(\boldsymbol{y}) = b^{-3}g(\boldsymbol{y}/b)$, $g = \abs{\psi}^2$ und $b = \massP/\massTot$, und zwar mit dem Gewicht, mit dem dieser Ort im Schwerpunktszustand vorkommt. Die Summe aller Kopien ist $\densE$. Die Rechnung dazu folgt.}
\end{tpcfigure}

Für den Produktzustand lässt sich $\densE$ ausrechnen. Wir halten den Elektronenort
$\posE = \boldsymbol{x}$ fest und integrieren über $\posP$. Statt über $\posP$
integrieren wir über $\posRel = \boldsymbol{x} - \posP$; dabei ist
$\dd^3\posP = \dd^3\posRel$, und nach \eqref{c11.eq.umkehr} liegt der Schwerpunkt bei
$\posCM = \boldsymbol{x} - (\massP/\massTot)\,\posRel$. Das ergibt
\begin{equation}\label{c11.eq.faltung}\claim{c11.cl.faltung}{derived}
  \densE(\boldsymbol{x})
  = \int \Bigl\lvert \Phi\Bigl(\boldsymbol{x} - \frac{\massP}{\massTot}\,\posRel\Bigr)
    \Bigr\rvert^2 \abs{\psi(\posRel)}^2 \,\dd^3\posRel .
\end{equation}
Genauso findet man für das Proton, mit $\posRel = \posE - \boldsymbol{x}$ und
$\posCM = \boldsymbol{x} + (\massE/\massTot)\,\posRel$,
\begin{equation}\label{c11.eq.protondichte}\claim{c11.cl.protondichte}{derived}
  \densP(\boldsymbol{x})
  = \int \Bigl\lvert \Phi\Bigl(\boldsymbol{x} + \frac{\massE}{\massTot}\,\posRel\Bigr)
    \Bigr\rvert^2 \abs{\psi(\posRel)}^2 \,\dd^3\posRel .
\end{equation}
Mit $f = \abs{\Phi}^2$, $g = \abs{\psi}^2$, $b = \massP/\massTot$ und
$g_b(\boldsymbol{y}) = b^{-3}g(\boldsymbol{y}/b)$ ist \eqref{c11.eq.faltung} die
Faltung $\densE = f * g_b$. Jeder mögliche Ort des Schwerpunkts trägt eine skalierte
Kopie der Relativdichte bei, gewichtet mit der Wahrscheinlichkeit dieses Orts. Die
Dichte des Elektrons im Labor ist diese skalierte Relativdichte, verschmiert mit der
Unschärfe des Schwerpunkts (Abbildung~\ref{c11.fig.faltung}).

\begin{tpcfigure}{c11.fig.faltung}
  \caption{Die Faltung \eqref{c11.eq.faltung} entlang einer Geraden durch die Mitte des
    Schwerpunktpakets, für $\Znuc = 1$. Dünn durchgezogen: die Relativdichte
    $\abs{\psinlm{1}{0}{0}}^2$ des Zustands 1s aus Abschnitt~\ref{c11.s5}. Gestrichelt:
    die Dichte $\abs{\Phi}^2$ eines Schwerpunktpakets in Form einer Glockenkurve; ihre
    Breite, $\num[c11.cl.wolke]{0{,}7}{}$ Längeneinheiten, ist für das Bild gewählt.
    Dick: die Dichte $\densE$ des Elektrons im Labor. Sie ist breiter und flacher als
    beide. Am markierten Punkt, eine Längeneinheit von der Mitte, ist
    $\densE\,\aRed^3 = \num[c11.cl.wolke]{0{,}0379}{}$. Längeneinheit ist der Bohrsche
    Radius $\aRed$ mit reduzierter Masse aus Abschnitt~\ref{c11.s4}.}
\end{tpcfigure}

In zwei Grenzfällen wird aus der verschmierten Dichte das gewohnte Bild.

\emph{Scharfer Schwerpunkt.} Wir ziehen das Paket $\abs{\Phi}^2$ immer enger um einen
Punkt $\posCM_0$ zusammen. Einen normierbaren Zustand der Breite null gibt es nicht;
gemeint ist eine Folge immer schmalerer Pakete. Im Grenzfall bleibt vom Integral
\eqref{c11.eq.faltung} nur der Beitrag übrig, bei dem das Argument von $\Phi$ gleich
$\posCM_0$ ist.

\emph{Unendlich schweres Proton.} Wächst die Protonenmasse über alle Grenzen, so geht
$\massP/\massTot$ gegen eins und $\massE/\massTot$ gegen null, und die reduzierte Masse
geht gegen $\massE$. Die Protonendichte \eqref{c11.eq.protondichte} wird dann zu
$\abs{\Phi}^2$: Das Proton sitzt im Schwerpunkt.

In Formeln, in der ersten Zeile für den scharfen Schwerpunkt bei $\posCM_0$, in der
zweiten, wenn zusätzlich $\massP \to \infty$ geht:
\begin{equation}\label{c11.eq.grenzfall}\claim{c11.cl.grenzfall}{derived}
  \begin{aligned}
    \densE(\boldsymbol{x})
    &\longrightarrow \Bigl(\frac{\massTot}{\massP}\Bigr)^{3}
      \Bigl\lvert \psi\Bigl(\frac{\massTot}{\massP}\,(\boldsymbol{x} - \posCM_0)\Bigr)
      \Bigr\rvert^2, \\
    \densE(\boldsymbol{x})
    &\longrightarrow \abs{\psi(\boldsymbol{x} - \posCM_0)}^2 .
  \end{aligned}
\end{equation}
Erst wenn beides zusammenkommt, ein scharf lokalisierter Schwerpunkt \emph{und} ein
unendlich schweres Proton, ist $\abs{\psi}^2$ die Dichte des Elektrons um einen festen
Punkt. Bei endlicher Protonenmasse ist die Elektronendichte um den Faktor
$\massP/\massTot$ gestaucht, und das Proton hat seinerseits eine kleine Wolke um den
Schwerpunkt.

\begin{merksatz}[Was ein Orbitalbild zeigt]
  Ein Orbitalbild zeigt $\abs{\psi(\posRel)}^2$, die Dichte der Relativkoordinate
  $\posRel = \posE - \posP$. Es zeigt nicht, wo sich das Elektron im Labor aufhält: Dazu
  muss man sagen, wie der Schwerpunkt präpariert ist, und die Faltung
  \eqref{c11.eq.faltung} ausführen.
\end{merksatz}

Auch bei der Energie muss man die beiden Hälften auseinanderhalten. Ist $\psi$ eine
Lösung von $\op{\HamRel}\psi = E\psi$, so ist der Erwartungswert der Gesamtenergie im
Produktzustand
\begin{equation}\label{c11.eq.energie}\claim{c11.cl.energie}{derived}
  \Etot = \frac{\hbar^2}{2\massTot} \int \abs{\nabla\Phi(\posCM)}^2 \,\dd^3\posCM + E ,
\end{equation}
sofern das Integral endlich ist. Der erste Summand ist die Bewegungsenergie des
Schwerpunkts; er hängt davon ab, wie das Atom als Ganzes präpariert wurde, und er ist
nie null, weil eine normierte Funktion $\Phi$ nicht konstant sein kann. Der zweite
Summand, $E$, ist die \emph{innere} Energie des Atoms. Die Energien $\En{\nqn}$, die wir
gleich berechnen, sind innere Energien. Ein ruhendes Atom mit scharfer Gesamtenergie ist
in diesem Modell kein normierbarer Zustand; der Schwerpunkt ist immer ein freies
Wellenpaket. Seine Wahrscheinlichkeit in jedem festen beschränkten Gebiet geht für
große Zeiten gegen null. Wenn zweite Momente existieren, wächst seine Breite für große
Zeiten; anfangs kann es schmaler werden\zoomref{c11.v1}.

\begin{tpcfigure}{c11.fig.energie}
  \caption{Innere Energie und Gesamtenergie nach \eqref{c11.eq.energie}, für $\Znuc = 1$ und in Einheiten der Rydberg-Energie $\RydRed$ aus Abschnitt~\ref{c11.s4}. Die dicken Striche sind die inneren Energien $\En{1}$, $\En{2}$, $\En{3}$ aus Theorem~\ref{c11.thm.spektrum}. Schraffiert ist jeweils der Bereich, in dem der Erwartungswert der Gesamtenergie liegen kann: Er beginnt bei der inneren Energie und hat kein oberes Ende, weil die Bewegungsenergie des Schwerpunkts jeden positiven Wert haben kann.}
\end{tpcfigure}

\begin{vertiefung}[id=c11.v1,
    titel={Warum der Schwerpunkt zerfließt},
    wozu={Der Kern behauptet, dass der Schwerpunkt eines Atoms kein Energieeigenzustand
      sein kann, sondern immer ein Wellenpaket ist, und dass dieses Paket zerfließt. Hier
      steht die Rechnung dazu.},
    voraussetzung={Kern~\ref{c11.s2} und \ref{c11.s3}; die Fouriertransformation aus
      Kapitel~\ref{c05}. Nicht nötig für die Abschnitte~\ref{c11.s4} bis \ref{c11.s7}.},
    zurueck={Die innere Dichte $\abs{\psi}^2$ eines Energieeigenzustands steht still;
      die Dichte des Schwerpunkts und mit ihr die Dichten von Elektron und Proton im
      Labor laufen auseinander.}]
  Der Schwerpunkt ist ein freies Teilchen der Masse $\massTot$. Wir schreiben
  $\FT{\Phi}(\boldsymbol{k})$ für die Fouriertransformierte von $\Phi(\posCM)$; sie ist
  wieder normiert, und aus $-\Lap_{\posCM}\Phi$ wird durch zweimaliges partielles
  Integrieren $\abs{\boldsymbol{k}}^2\,\FT{\Phi}$.

  \begin{proposition}[Kein normierbarer Eigenzustand]\label{c11.prop.kein-eigenzustand}\claim[proof-included]{c11.cl.kein-eigenzustand}{derived}
    Es gibt keine reelle Zahl $\lambda$ und keine quadratintegrierbare Funktion
    $\Phi \neq 0$ mit $-\Lap_{\posCM}\Phi = \lambda\Phi$.
  \end{proposition}
  \begin{proof}
    Nach Fouriertransformation lautet die Gleichung
    $(\abs{\boldsymbol{k}}^2 - \lambda)\,\FT{\Phi}(\boldsymbol{k}) = 0$ für fast alle
    $\boldsymbol{k}$. Für $\lambda < 0$ ist der Vorfaktor nirgends null. Für
    $\lambda \geq 0$ ist er nur auf einer Kugelfläche (für $\lambda = 0$: in einem Punkt)
    null, und das ist eine Menge vom Volumen null. In beiden Fällen ist $\FT{\Phi}$ fast
    überall null, also $\Phi = 0$.
  \end{proof}

  Eine ebene Welle $e^{i\boldsymbol{k}\cdot\posCM}$ löst die Eigenwertgleichung zwar,
  hat aber überall den Betrag eins und ist nicht normierbar. Ein normierbarer
  Schwerpunktszustand ist deshalb immer eine Überlagerung vieler Impulse, ein
  Wellenpaket. Wie breit es ist, misst die Varianz $\Var\posCM$, die Summe der Varianzen
  der drei Komponenten. Für ein freies Teilchen ist im Heisenbergbild
  $\posCM(t) = \posCM(0) + t\,\boldsymbol{P}/\massTot$ mit dem Gesamtimpuls
  $\boldsymbol{P}$. Mit den Abweichungen $\delta\posCM$ und $\delta\boldsymbol{P}$ von
  den Erwartungswerten folgt für Pakete, deren zweite Momente endlich sind,
  \begin{equation}\label{c11.eq.varianz}\claim{c11.cl.varianz}{derived}
    \Var\posCM(t) = \Var\posCM(0)
    + \frac{t}{\massTot}\,\expval{\delta\posCM \cdot \delta\boldsymbol{P}
      + \delta\boldsymbol{P} \cdot \delta\posCM}
    + \frac{t^2}{\massTot^2}\,\Var\boldsymbol{P}.
  \end{equation}
  Die Impulsvarianz $\Var\boldsymbol{P}$ eines normierbaren Pakets ist nicht null. Für
  große Zeiten gewinnt also der letzte Term, und die Breite wächst proportional zu
  $\abs{t}$. Am Anfang kann das Paket schmaler werden, wenn der mittlere Term negativ
  ist (Abbildung~\ref{c11.fig.zerfliessen}).

  \begin{tpcfigure}{c11.fig.zerfliessen}
    \caption{Die Breite eines freien Schwerpunktpakets nach \eqref{c11.eq.varianz},
      bezogen auf die Anfangsbreite. Waagerecht die Zeit, gemessen in Einheiten von
      $\massTot\sqrt{\Var\posCM(0)}/\sqrt{\Var\boldsymbol{P}}$. Die Zahl $k$ ist der
      Erwartungswert im mittleren Term von \eqref{c11.eq.varianz}, geteilt durch
      $2\sqrt{\Var\posCM(0)\,\Var\boldsymbol{P}}$. Durchgezogen: $k = 0$. Gestrichelt:
      ein Paket mit $k < 0$, das erst schmaler und dann breiter wird; der Punkt markiert
      seine kleinste Breite. Dünn: eine Gerade mit der gemeinsamen asymptotischen
      Steigung.}
  \end{tpcfigure}

  Ohne Annahme über zweite Momente gilt eine schwächere, aber allgemeine Aussage.

  \begin{proposition}[Lokales Zerfließen]\label{c11.prop.zerfliessen}\claim[proof-included]{c11.cl.zerfliessen}{derived}
    Sei $\Phi_t$ die freie Zeitentwicklung eines normierten $\Phi_0$ und $B$ ein
    beschränktes Gebiet. Dann geht die Wahrscheinlichkeit
    $\int_B \abs{\Phi_t(\posCM)}^2 \,\dd^3\posCM$, den Schwerpunkt in $B$ zu finden, für
    $\abs{t} \to \infty$ gegen null.
  \end{proposition}
  \begin{proof}
    Die freie Zeitentwicklung multipliziert $\FT{\Phi}(\boldsymbol{k})$ mit
    $e^{-i\hbar\abs{\boldsymbol{k}}^2 t/(2\massTot)}$. Zurück im Ortsraum ist das ein
    Integral $\Phi_t(\posCM) = \int K_t(\posCM - \posCM')\,\Phi_0(\posCM')\,\dd^3\posCM'$
    mit einem Kern, dessen Betrag nicht vom Ort abhängt:
    \begin{equation}\label{c11.eq.freikern}
      K_t(\boldsymbol{z}) = \Bigl(\frac{\massTot}{2\pi i\hbar t}\Bigr)^{3/2}
        e^{i\massTot\abs{\boldsymbol{z}}^2/(2\hbar t)},
      \qquad
      \abs{K_t(\boldsymbol{z})} = \Bigl(\frac{\massTot}{2\pi\hbar\abs{t}}\Bigr)^{3/2}.
    \end{equation}
    Der Kern entsteht aus einem gaußschen Integral über $\boldsymbol{k}$. Für ein
    glattes, rasch abfallendes $g$ anstelle von $\Phi_0$ ist deshalb überall
    \begin{equation}\label{c11.eq.schranke}
      \abs{g_t(\posCM)} \leq \Bigl(\frac{\massTot}{2\pi\hbar\abs{t}}\Bigr)^{3/2}
        \int \abs{g(\posCM')}\,\dd^3\posCM' ,
    \end{equation}
    und das Integral von $\abs{g_t}^2$ über $B$ ist höchstens das Volumen von $B$ mal
    dem Quadrat dieser Schranke; es geht gegen null. Ein beliebiges normiertes $\Phi_0$
    lässt sich durch solche $g$ im quadratischen Mittel beliebig genau annähern, und die
    freie Zeitentwicklung ändert den Abstand im quadratischen Mittel nicht. Also liegt
    die Wurzel der gesuchten Wahrscheinlichkeit für große $\abs{t}$ unter jeder
    vorgegebenen Schranke.
  \end{proof}

  Die Aussage sagt nicht, dass die Dichte an jedem Punkt abnimmt, und nicht, dass die
  Breite von Anfang an wächst. Für die Relativbewegung ist es anders: Ein
  Energieeigenzustand entwickelt sich als
  $\psi(t) = e^{-i\En{\nqn}t/\hbar}\,\psi(0)$, seine Dichte $\abs{\psi}^2$ ändert sich
  nicht.
\end{vertiefung}


\section{Die Relativbewegung: Radialgleichung und Energien}\label{c11.s4}

Wir suchen jetzt die Lösungen von
$\op{\HamRel}\psi = E\psi$ mit negativer Energie $E$, die gebundenen Zustände. Das ist das
Problem eines einzelnen Teilchens der Masse $\massRed$ in einem Potential, das nur vom
Abstand $\rad = \abs{\posRel}$ abhängt. Für solche Potentiale haben wir in
Kapitel~\ref{c10} Kugelkoordinaten $(\rad, \theta, \phi)$ benutzt und den Winkelanteil
ein für alle Mal gelöst: mit den Kugelflächenfunktionen $\Ylm{\lqn}{\mqn}$ aus
\defref{D67}. Von dort übernehmen wir eine einzige Tatsache. Auf einem Produkt aus einer
Funktion $R(\rad)$ des Abstands und einer Kugelflächenfunktion wirkt der
Laplace-Operator so:
\begin{equation}\label{c11.eq.winkel}\claim{c11.cl.winkel}{numerically-checked}
  \Lap_{\posRel}\bigl(R(\rad)\,\Ylm{\lqn}{\mqn}(\theta, \phi)\bigr)
  = \Bigl(\frac{1}{\rad^2}\frac{\dd}{\dd\rad}\Bigl(\rad^2\frac{\dd R}{\dd\rad}\Bigr)
    - \frac{\lqn(\lqn + 1)}{\rad^2}\,R\Bigr)\,\Ylm{\lqn}{\mqn}(\theta, \phi).
\end{equation}
Die Winkel sind damit erledigt; die Zahl $\lqn$ merkt sich, wie viel Drehimpuls in der
Bewegung steckt. Für $\lqn = 0$, $1$ und $2$, mehr brauchen wir in diesem Kapitel nicht,
ist \eqref{c11.eq.winkel} für beliebiges $R$ exakt nachgerechnet.

Es bleibt eine Gleichung für $R(\rad)$. Sie wird noch einfacher, wenn man statt $R$ die
Funktion $u(\rad) = \rad\,R(\rad)$ benutzt, denn dann ist
\begin{equation}\label{c11.eq.usubst}\claim{c11.cl.usubst}{derived}
  \frac{1}{\rad^2}\frac{\dd}{\dd\rad}\Bigl(\rad^2\frac{\dd}{\dd\rad}\frac{u}{\rad}\Bigr)
  = \frac{1}{\rad}\frac{\dd^2 u}{\dd\rad^2} .
\end{equation}
Den Rest fassen wir zu einem Potential zusammen.

\tpcdef{D71}

Setzt man $\psi = (u/\rad)\,\Ylm{\lqn}{\mqn}$ in $\op{\HamRel}\psi = E\psi$ ein, benutzt
\eqref{c11.eq.winkel} und multipliziert mit $\rad$, so erhält man die
\emph{Radialgleichung}
\begin{equation}\label{c11.eq.radial}\claim{c11.cl.radial}{derived}
  \begin{aligned}
    -\frac{\hbar^2}{2\massRed}\,\frac{\dd^2 u}{\dd\rad^2} + \Veff{\lqn}(\rad)\,u(\rad)
    &= E\,u(\rad), \\
    \Veff{\lqn}(\rad) &= \frac{\hbar^2\lqn(\lqn + 1)}{2\massRed\,\rad^2}
      - \frac{\Znuc\echarge^2}{4\pi\epsz\,\rad}.
  \end{aligned}
\end{equation}
Das ist die Schrödingergleichung eines Teilchens auf der Halbgeraden $\rad > 0$ im
effektiven Potential $\Veff{\lqn}$, also ein eindimensionales Problem.
Abbildung~\ref{c11.fig.potential} zeigt dieses Potential. Wir suchen Lösungen, die
normierbar sind. Weil die
Kugelflächenfunktionen auf der Kugel normiert sind, heißt das
\begin{equation}\label{c11.eq.unorm}\claim{c11.cl.unorm}{derived}
  \int_0^\infty \abs{R(\rad)}^2\,\rad^2\,\dd\rad = \int_0^\infty \abs{u(\rad)}^2\,\dd\rad = 1 .
\end{equation}
Der Faktor $\rad^2$ kommt vom Volumenelement in Kugelkoordinaten; er ist der Grund,
warum $u$ bequemer ist als $R$.

\begin{tpcfigure}{c11.fig.potential}
  \caption{Das effektive Potential $\Veff{\lqn}$ aus \eqref{c11.eq.radial} für
    $\Znuc = 1$. Durchgezogen: $\lqn = 0$, das reine Coulombpotential. Gestrichelt:
    $\lqn = 1$; der Zentrifugalterm hält das Teilchen vom Ursprung fern und erzeugt eine
    Mulde, deren tiefster Punkt markiert ist. Dünne waagerechte Linien: die Energien
    $\En{1}$, $\En{2}$, $\En{3}$ aus Theorem~\ref{c11.thm.spektrum}. Längen in Einheiten
    von $\aRed$, Energien in Einheiten von $\RydRed$; beide enthalten die reduzierte
    Masse.}
\end{tpcfigure}

\subsection{Die natürlichen Einheiten des Atoms}\label{c11.s4.1}

Die Konstanten in \eqref{c11.eq.radial} legen eine Länge und eine Energie fest. In
diesen Einheiten wird die Gleichung übersichtlich.

\tpcdef{D63}

\tpcdef{D64}

Es gibt also zwei Bohrsche Radien, und man muss bei jeder Formel sagen, welchen man
meint. In die Radialgleichung geht die reduzierte Masse ein, also $\aRed$. Aus den
beiden Definitionen folgt
\begin{equation}\label{c11.eq.bohr}\claim{c11.cl.bohr}{derived}
  \begin{aligned}
    \aRed &= \frac{4\pi\epsz\hbar^2}{\massRed\,\echarge^2}
           = \aBohr\,\frac{\massE}{\massRed}
           = \aBohr\,\Bigl(1 + \frac{\massE}{\massP}\Bigr), \\
    \RydRed &= \frac{\hbar^2}{2\massRed\,\aRed^2}
             = \frac{\echarge^2}{8\pi\epsz\,\aRed}
             = \frac{\massRed}{\massE}\,\Ryd .
  \end{aligned}
\end{equation}
Die Energie $\RydRed$ ist die Rydberg-Energie mit reduzierter Masse. In Zahlen:
$\aBohr = \num[c11.cl.a0]{5{,}29177 \cdot 10^{-11}}{m}$ und
$\aRed = \num[c11.cl.amu]{5{,}29465 \cdot 10^{-11}}{m}$. Weiter ist
$\Ryd = \num[c11.cl.ry]{13{,}6057}{eV}$ und
$\RydRed = \num[c11.cl.e1-mu]{13{,}5983}{eV}$.

Mit der dimensionslosen Länge $x = \Znuc\rad/\aRed$ und der dimensionslosen Energie
$\varepsilon = E/(\Znuc^2\,\RydRed)$ verschwinden alle Konstanten aus der
Radialgleichung. Aus \eqref{c11.eq.radial} wird, mit $u$ als Funktion von $x$,
\begin{equation}\label{c11.eq.dimlos}\claim{c11.cl.dimlos}{derived}
  -\frac{\dd^2 u}{\dd x^2} + \Bigl(\frac{\lqn(\lqn + 1)}{x^2} - \frac{2}{x}\Bigr)\,u
  = \varepsilon\,u .
\end{equation}

\subsection{Lösungen durch Einsetzen}\label{c11.s4.2}

Gleichungen wie \eqref{c11.eq.dimlos} löst man systematisch mit einem Potenzreihenansatz.
Das ist aufwendig. Wir gehen den kurzen Weg: Wir schreiben Lösungen hin und prüfen sie
durch Einsetzen. Die einfachste ist $u(x) = x\,e^{-x}$ für $\lqn = 0$. Ihre Ableitungen
sind $u' = (1 - x)\,e^{-x}$ und $u'' = (x - 2)\,e^{-x}$, und damit
\begin{equation}\label{c11.eq.grundzustand}\claim{c11.cl.grundzustand}{derived}
  -u'' - \frac{2}{x}\,u = (2 - x)\,e^{-x} - 2\,e^{-x} = -x\,e^{-x} = -u .
\end{equation}
Die Funktion löst \eqref{c11.eq.dimlos} also mit $\varepsilon = -1$, das heißt mit der
Energie $E = -\Znuc^2\,\RydRed$. Sie verschwindet bei $x = 0$ und fällt für große $x$
exponentiell ab; sie ist normierbar. Sie hat die niedrigste Energie in der
konstruierten Familie. Auf dieselbe Weise prüft man jede der folgenden Lösungen.

\begin{tpcfigure}{c11.fig.grundzustand}
  \caption{Die Lösung $u(x) = x\,e^{-x}$ aus \eqref{c11.eq.grundzustand} im Potential, für $\Znuc = 1$. Dick: das effektive Potential für $\lqn = 0$. Dünne waagerechte Linie: die Energie $\En{1} = -\RydRed$. Die Funktion $u$ ist, wie in solchen Bildern üblich, auf dieser Linie als Grundlinie gezeichnet; ihre Höhe ist frei skaliert. Sie beginnt bei null, hat ihr Maximum bei $\rad = \aRed$ und fällt nach außen exponentiell ab.}
\end{tpcfigure}

\begin{theorem}[Energien der konstruierten gebundenen Zustände]\label{c11.thm.spektrum}\claim[proof-linked]{c11.cl.spektrum}{derived}
  Sei $\nqn$ eine der Zahlen $1, 2, 3, \ldots$ und $\lqn$ eine der Zahlen
  $0, 1, \ldots, \nqn - 1$. Dann hat die Radialgleichung \eqref{c11.eq.radial} eine
  normierbare Lösung $u \neq 0$ mit $u(0) = 0$ zur Energie
  \begin{equation}\label{c11.eq.energien}
    \En{\nqn} = -\frac{\Znuc^2\,\RydRed}{\nqn^2}
              = -\frac{\Znuc^2\echarge^2}{8\pi\epsz\,\aRed}\,\frac{1}{\nqn^2} .
  \end{equation}
  Die Energie hängt nur von $\nqn$ ab, nicht von $\lqn$.
\end{theorem}

Für $\nqn = 1$ haben wir das eben nachgerechnet; die Lösungen für $\nqn = 2$ stehen in
Abschnitt~\ref{c11.s5}, und Aufgabe~\ref{c11.exc.zweip} prüft eine davon. Für größere
$\nqn$ geht man genauso vor; den allgemeinen Beweis führen wir hier nicht\zoomref{c11.v2}.
Der Satz sagt, dass es diese Lösungen gibt. Dass die Relativbewegung keine weiteren gebundenen Zustände hat, zeigen
wir in diesem Kapitel nicht.

\tpcdef{D65}

Die Zählung der konstruierten Familie beginnt bei $\nqn = 1$. Ihre niedrigste Energie
ist $\En{1} = -\Znuc^2\,\RydRed$, die des Zustands 1s. Darüber liegen in dieser Familie
unendlich viele weitere Niveaus, die immer dichter zusammenrücken und sich bei $E = 0$
häufen (Abbildung~\ref{c11.fig.termschema} in Abschnitt~\ref{c11.s6}). Für Wasserstoff ist
$\En{1} = \num[c11.cl.e1-mu]{-13{,}5983}{eV}$; mit der Elektronenmasse anstelle der
reduzierten Masse käme $\num[c11.cl.e1-me]{-13{,}6057}{eV}$ heraus.

Zu jedem $\nqn$ gehören mehrere Zustände gleicher Energie. Gezählt werden sie mit
\defref{D65}: Zu jedem $\lqn$ gibt es $2\lqn + 1$ Werte von $\mqn$, also insgesamt
\begin{equation}\label{c11.eq.entartung}\claim{c11.cl.entartung}{derived}
  \sum_{\lqn = 0}^{\nqn - 1} (2\lqn + 1) = \nqn^2 .
\end{equation}

\begin{tpctable}{c11.tab.zustaende}
  \caption{Die Zustände der Relativbewegung bis $\nqn = 3$. Die Buchstaben s, p, d
    stehen für $\lqn = 0, 1, 2$; in der letzten Spalte steht die Zahl der Zustände
    gleicher Energie nach \eqref{c11.eq.entartung}. Der Spin ist nicht mitgezählt.}
  \begin{tabular}{ccccc}
    \toprule
    $\nqn$ & $\lqn$ & $\mqn$ & Name & Zustände zu $\En{\nqn}$ \\
    \midrule
    $1$ & $0$ & $0$ & 1s & $1$ \\
    \midrule
    $2$ & $0$ & $0$ & 2s & \\
    $2$ & $1$ & $-1, 0, 1$ & 2p & $4$ \\
    \midrule
    $3$ & $0$ & $0$ & 3s & \\
    $3$ & $1$ & $-1, 0, 1$ & 3p & \\
    $3$ & $2$ & $-2, -1, 0, 1, 2$ & 3d & $9$ \\
    \bottomrule
  \end{tabular}
\end{tpctable}

\begin{tpcfigure}{c11.fig.entartung}
  \caption{Dieselben Zustände wie in Tabelle~\ref{c11.tab.zustaende}, nach $\nqn$ und $\lqn$ sortiert, für $\Znuc = 1$. Jeder kurze Strich ist ein Zustand; seine Höhe ist die Energie $\En{\nqn}$ in Einheiten von $\RydRed$. Striche auf gleicher Höhe haben dieselbe Energie, gleichgültig in welcher Spalte sie stehen. Rechts steht ihre Anzahl $\nqn^2$.}
\end{tpcfigure}

\begin{vertiefung}[id=c11.v2,
    titel={Die allgemeine Radialfunktion},
    wozu={Der Kern prüft die Lösungen für $\nqn \leq 2$ durch Einsetzen. Hier steht die
      Formel für alle $\nqn$ und $\lqn$, der Beweis von Theorem~\ref{c11.thm.spektrum}
      und der Beweis, dass der Vorfaktor die Radialfunktionen normiert.},
    voraussetzung={Kern~\ref{c11.s4}. Nicht nötig für die Abschnitte~\ref{c11.s5} bis
      \ref{c11.s7}, die mit den ausgeschriebenen Funktionen für $\nqn \leq 2$ arbeiten.},
    zurueck={Zu jedem $\nqn$ und jedem $\lqn < \nqn$ gibt es eine normierte Radialfunktion
      zur Energie $\En{\nqn}$: ein Polynom vom Grad $\nqn - \lqn - 1$ mal
      $\rad^{\lqn}$ mal einer abfallenden Exponentialfunktion.}]
  Die Polynome, die in den Radialfunktionen auftreten, haben einen Namen. Wie man sie
  normiert und beschriftet, ist eine Konvention; wir legen sie fest.

  \tpcdef{D68}

  Die ersten drei lauten
  \begin{align}
    \Lag{0}{\alpha}(x) &= 1,
      \label{c11.eq.lag0}\claim{c11.cl.lag0}{numerically-checked} \\
    \Lag{1}{\alpha}(x) &= \alpha + 1 - x,
      \label{c11.eq.lag1}\claim{c11.cl.lag1}{numerically-checked} \\
    \Lag{2}{\alpha}(x) &= \frac{x^2}{2} - (\alpha + 2)\,x + \frac{(\alpha + 1)(\alpha + 2)}{2}.
      \label{c11.eq.lag2}\claim{c11.cl.lag2}{numerically-checked}
  \end{align}

  \begin{tpcfigure}{c11.fig.laguerre}
    \caption{Die Polynome $\Lag{k}{1}$ für $k = 0, 1, 2, 3$ nach \defref{D68}. Das Polynom zum Index $k$ hat den Grad $k$ und bei $x = 0$ den Wert $k + 1$. Mit $\alpha = 1$ sind das die Polynome, die in den Radialfunktionen der s-Zustände stehen.}
  \end{tpcfigure}

  \tpcdef{D66}

  \begin{proof}[Beweis von Theorem~\ref{c11.thm.spektrum}]
    Für $E < 0$ setzen wir $\kappa = \sqrt{-2\massRed E}/\hbar$, $s = 2\kappa\rad$ und
    $u = s^{\lqn + 1}\,e^{-s/2}\,v(s)$. Einsetzen in \eqref{c11.eq.radial} und Kürzen
    ergibt für $v$ die Gleichung
    \begin{equation}\label{c11.eq.laguerre-dgl}\claim{c11.cl.laguerre-dgl}{derived}
      s\,v'' + (2\lqn + 2 - s)\,v' + k\,v = 0,
      \qquad
      k = \frac{\Znuc}{\kappa\,\aRed} - \lqn - 1 .
    \end{equation}
    Mit dem Potenzreihenansatz $v = \sum_j c_j\,s^j$ folgt durch Koeffizientenvergleich
    \begin{equation}\label{c11.eq.rekursion}\claim{c11.cl.rekursion}{derived}
      c_{j + 1} = \frac{j - k}{(j + 1)(j + 2\lqn + 2)}\;c_j .
    \end{equation}
    Ist $k$ eine der Zahlen $0, 1, 2, \ldots$, so ist $c_{k + 1} = 0$ und mit ihm alle
    folgenden Koeffizienten: $v$ ist ein Polynom vom Grad $k$. Wir schreiben
    $k = \nqn - \lqn - 1$; dann ist $\nqn \geq \lqn + 1$ eine ganze Zahl, und aus der
    Definition von $k$ folgt
    \begin{equation}\label{c11.eq.abbruch}\claim{c11.cl.abbruch}{derived}
      \kappa = \frac{\Znuc}{\nqn\,\aRed},
      \qquad
      E = -\frac{\hbar^2\kappa^2}{2\massRed}
        = -\frac{\hbar^2\Znuc^2}{2\massRed\,\aRed^2\,\nqn^2}
        = -\frac{\Znuc^2\,\RydRed}{\nqn^2} .
    \end{equation}
    Das Polynom ist bis auf einen Faktor $\Lag{\nqn - \lqn - 1}{2\lqn + 1}(s)$: Das
    Verhältnis aufeinanderfolgender Koeffizienten der Summe in \defref{D68} ist für
    $\alpha = 2\lqn + 1$ gerade das in \eqref{c11.eq.rekursion}. Die Lösung
    $u = s^{\lqn + 1}\,e^{-s/2}\,v$ verschwindet bei $\rad = 0$ und fällt exponentiell
    ab, ist also normierbar, und $R = u/\rad$ ist bis auf den Vorfaktor die
    Radialfunktion aus \defref{D66}.
  \end{proof}

  Der Beweis konstruiert Lösungen. Er zeigt nicht, dass eine Potenzreihe, die nicht
  abbricht, niemals normierbar ist; diese Frage bleibt hier offen.

  \begin{proposition}[Normierung der Radialfunktionen]\label{c11.prop.norm}\claim[proof-included]{c11.cl.radialnorm}{derived}
    Für alle $\nqn \geq 1$ und $0 \leq \lqn \leq \nqn - 1$ gilt
    $\int_0^\infty \Rnl{\nqn}{\lqn}(\rad)^2\,\rad^2\,\dd\rad = 1$.
  \end{proposition}
  \begin{proof}
    Wir schreiben $k = \nqn - \lqn - 1$ und $\alpha = 2\lqn + 1$. Aus der Summe in
    \defref{D68} folgt durch Ausdifferenzieren mit der Produktregel die Darstellung
    \begin{equation}\label{c11.eq.rodrigues}
      \Lag{k}{\alpha}(s) = \frac{s^{-\alpha}\,e^{s}}{k!}\,
        \frac{\dd^k}{\dd s^k}\bigl(e^{-s}\,s^{k + \alpha}\bigr).
    \end{equation}
    Integriert man $s^\alpha e^{-s}\,\Lag{k}{\alpha}$ gegen ein Polynom und wälzt die
    $k$ Ableitungen durch partielle Integration auf das Polynom, so verschwindet das
    Integral für jedes Polynom vom Grad kleiner als $k$; die Randterme sind null. Gegen
    $\Lag{k}{\alpha}$ selbst, dessen höchster Koeffizient $(-1)^k/k!$ ist, bleibt
    \begin{equation}\label{c11.eq.lagnorm}
      \int_0^\infty s^{\alpha}\,e^{-s}\,\bigl(\Lag{k}{\alpha}(s)\bigr)^2\,\dd s
      = \frac{(k + \alpha)!}{k!} .
    \end{equation}
    Für die Radialfunktionen brauchen wir eine Potenz von $s$ mehr. Koeffizientenvergleich
    in den Summen ergibt, mit $\Lag{-1}{\alpha} = 0$,
    \begin{equation}\label{c11.eq.dreiterm}
      s\,\Lag{k}{\alpha} = (2k + \alpha + 1)\,\Lag{k}{\alpha}
        - (k + 1)\,\Lag{k + 1}{\alpha} - (k + \alpha)\,\Lag{k - 1}{\alpha} .
    \end{equation}
    Multipliziert mit $s^\alpha e^{-s}\,\Lag{k}{\alpha}$ und integriert fallen die beiden
    letzten Terme wegen der Orthogonalität weg, und mit \eqref{c11.eq.lagnorm} bleibt
    \begin{equation}\label{c11.eq.lagnorm2}
      \int_0^\infty s^{\alpha + 1}\,e^{-s}\,\bigl(\Lag{k}{\alpha}(s)\bigr)^2\,\dd s
      = (2k + \alpha + 1)\,\frac{(k + \alpha)!}{k!}
      = \frac{2\nqn\,(\nqn + \lqn)!}{(\nqn - \lqn - 1)!} .
    \end{equation}
    Mit $s = 2\kappa\rad$ ist $\rad^2\,\dd\rad = s^2\,\dd s/(2\kappa)^3$, und in
    $\Rnl{\nqn}{\lqn}^2\,\rad^2$ steht $s^{2\lqn + 2} = s^{\alpha + 1}$. Das
    Normierungsintegral ist also das Quadrat des Vorfaktors aus \defref{D66}, geteilt
    durch $(2\kappa)^3$, mal der rechten Seite von \eqref{c11.eq.lagnorm2}, und das ist
    eins.
  \end{proof}
\end{vertiefung}

\section{Die ersten Wellenfunktionen}\label{c11.s5}


Jetzt setzen wir die Teile zusammen und schreiben die Wellenfunktionen der
Relativbewegung für $\nqn = 1$ und $\nqn = 2$ vollständig hin. Dazu brauchen wir die
Kugelflächenfunktionen für $\lqn = 0$ und $\lqn = 1$. Nach \defref{D67} sind das
\begin{align}
  \Ylm{0}{0}(\theta, \phi) &= \frac{1}{\sqrt{4\pi}},
    \label{c11.eq.y00}\claim{c11.cl.y00}{numerically-checked} \\
  \Ylm{1}{0}(\theta, \phi) &= \sqrt{\frac{3}{4\pi}}\,\cos\theta,
    \label{c11.eq.y10}\claim{c11.cl.y10}{numerically-checked} \\
  \Ylm{1}{\pm 1}(\theta, \phi) &= \mp\sqrt{\frac{3}{8\pi}}\,\sin\theta\,e^{\pm i\phi}.
    \label{c11.eq.y1pm}\claim{c11.cl.y1pm}{numerically-checked}
\end{align}
In der letzten Zeile gehören die oberen und die unteren Vorzeichen jeweils zusammen. Das
Minuszeichen vor $\Ylm{1}{1}$ ist kein Naturgesetz, sondern eine Verabredung, die
Condon-Shortley-Konvention. Jede andere Wahl der Vorzeichen beschriebe dieselbe Physik;
man muss nur bei einer bleiben.

Zu den Lösungen $u$ aus Theorem~\ref{c11.thm.spektrum} gehören die Radialfunktionen
$\Rnl{\nqn}{\lqn} = u/\rad$. Normiert lauten sie
\begin{align}
  \Rnl{1}{0}(\rad) &= 2\left(\frac{\Znuc}{\aRed}\right)^{3/2} e^{-\Znuc\rad/\aRed},
    \label{c11.eq.r10}\claim{c11.cl.r10}{numerically-checked} \\
  \Rnl{2}{0}(\rad) &= \frac{1}{2\sqrt{2}}\left(\frac{\Znuc}{\aRed}\right)^{3/2}
    \left(2 - \frac{\Znuc\rad}{\aRed}\right) e^{-\Znuc\rad/(2\aRed)},
    \label{c11.eq.r20}\claim{c11.cl.r20}{numerically-checked} \\
  \Rnl{2}{1}(\rad) &= \frac{1}{2\sqrt{6}}\left(\frac{\Znuc}{\aRed}\right)^{3/2}
    \frac{\Znuc\rad}{\aRed}\,e^{-\Znuc\rad/(2\aRed)}.
    \label{c11.eq.r21}\claim{c11.cl.r21}{numerically-checked}
\end{align}
In allen dreien steht der Bohrsche Radius $\aRed$ mit der reduzierten Masse
(Abbildung~\ref{c11.fig.radial}).

\begin{tpcfigure}{c11.fig.radial}
  \caption{Die Radialfunktionen \eqref{c11.eq.r10} bis \eqref{c11.eq.r21} für
    $\Znuc = 1$. $\Rnl{1}{0}$ und $\Rnl{2}{0}$ sind am Ursprung von null verschieden,
    $\Rnl{2}{1}$ beginnt bei null. $\Rnl{2}{0}$ wechselt bei $\rad = 2\aRed$ das
    Vorzeichen; eine solche Nullstelle heißt Knoten.}
\end{tpcfigure}

\tpcdef{D69}

Ausmultipliziert ergibt das für $\nqn = 1$ und $\nqn = 2$ fünf Funktionen:
\begin{align}
  \psinlm{1}{0}{0}(\rad, \theta, \phi) &= \frac{1}{\sqrt{\pi}}
    \left(\frac{\Znuc}{\aRed}\right)^{3/2} e^{-\Znuc\rad/\aRed},
    \label{c11.eq.psi100}\claim{c11.cl.psi100}{numerically-checked} \\
  \psinlm{2}{0}{0}(\rad, \theta, \phi) &= \frac{1}{4\sqrt{2\pi}}
    \left(\frac{\Znuc}{\aRed}\right)^{3/2}
    \left(2 - \frac{\Znuc\rad}{\aRed}\right) e^{-\Znuc\rad/(2\aRed)},
    \label{c11.eq.psi200}\claim{c11.cl.psi200}{numerically-checked} \\
  \psinlm{2}{1}{0}(\rad, \theta, \phi) &= \frac{1}{4\sqrt{2\pi}}
    \left(\frac{\Znuc}{\aRed}\right)^{3/2}
    \frac{\Znuc\rad}{\aRed}\,e^{-\Znuc\rad/(2\aRed)}\,\cos\theta,
    \label{c11.eq.psi210}\claim{c11.cl.psi210}{numerically-checked} \\
  \psinlm{2}{1}{\pm 1}(\rad, \theta, \phi) &= \mp\frac{1}{8\sqrt{\pi}}
    \left(\frac{\Znuc}{\aRed}\right)^{3/2}
    \frac{\Znuc\rad}{\aRed}\,e^{-\Znuc\rad/(2\aRed)}\,\sin\theta\,e^{\pm i\phi}.
    \label{c11.eq.psi21pm}\claim{c11.cl.psi21pm}{numerically-checked}
\end{align}
Die erste heißt 1s, die zweite 2s, die drei übrigen 2p. Die Zustände 1s und 2s hängen
nicht von den Winkeln ab; sie sind kugelsymmetrisch. Der Zustand 1s ist am Ursprung
am größten und fällt nach außen exponentiell ab. Bei den 2p-Zuständen kommt die
Winkelabhängigkeit dazu (Abbildung~\ref{c11.fig.schnitt}).

\begin{tpcfigure}{c11.fig.schnitt}
  \demo{c11.demo.cloud}
  \caption{Die Dichte $\abs{\psi}^2$ der Relativkoordinate für die 2p-Zustände, als
    Schnitt durch die Ebene $y = 0$ und für $\Znuc = 1$; dunkler heißt größer, beide
    Felder benutzen dieselbe lineare Grauskala. Links $\psinlm{2}{1}{0}$: zwei Keulen
    entlang der $z$-Achse, getrennt durch die Ebene $z = 0$, in der die Dichte
    verschwindet. Rechts $\psinlm{2}{1}{\pm 1}$: Im Raum ist das ein Ring um die
    $z$-Achse, auf der die Dichte verschwindet; der Schnitt trifft den Ring zweimal. Die
    Achsen sind Komponenten von $\posRel = \posE - \posP$, nicht Orte im Labor.}
\end{tpcfigure}

Dass diese Funktionen normiert sind, rechnen wir für den Zustand 1s vor. In
Kugelkoordinaten ist das Volumenelement $\rad^2\sin\theta\,\dd\rad\,\dd\theta\,\dd\phi$;
weil $\psinlm{1}{0}{0}$ nicht von den Winkeln abhängt, liefert das Winkelintegral die
Kugeloberfläche $4\pi$, und mit $x = \Znuc\rad/\aRed$ bleibt
\begin{equation}\label{c11.eq.norm100}\claim{c11.cl.norm100}{derived}
  \begin{aligned}
    \int \abs{\psinlm{1}{0}{0}}^2\,\dd^3\posRel
    &= \frac{1}{\pi}\Bigl(\frac{\Znuc}{\aRed}\Bigr)^{3}\,4\pi
      \int_0^\infty \rad^2\,e^{-2\Znuc\rad/\aRed}\,\dd\rad \\
    &= 4\int_0^\infty x^2\,e^{-2x}\,\dd x
     = 4 \cdot \frac{2}{2^3} = 1 .
  \end{aligned}
\end{equation}
Benutzt haben wir $\int_0^\infty x^k\,e^{-cx}\,\dd x = k!/c^{k + 1}$, hier mit $k = 2$
und $c = 2$. Für die anderen vier Funktionen geht es genauso, nur mit mehr Schreibarbeit.
Alle fünf sind normiert, stehen paarweise senkrecht aufeinander und lösen die Gleichung
der Relativbewegung zu der Energie aus Theorem~\ref{c11.thm.spektrum}. Die
Eigenfunktionsgleichung gilt für $\posRel \neq 0$ und als $L^2$-Identität fast überall:
\begin{equation}\label{c11.eq.orthonormal}\claim{c11.cl.orthonormal}{numerically-checked}
  \begin{aligned}
    \int \overline{\psinlm{\nqn}{\lqn}{\mqn}}\;\psinlm{\nqn'}{\lqn'}{\mqn'}\,\dd^3\posRel
    &= \delta_{\nqn\nqn'}\,\delta_{\lqn\lqn'}\,\delta_{\mqn\mqn'}, \\
    \op{\HamRel}\,\psinlm{\nqn}{\lqn}{\mqn} &= \En{\nqn}\,\psinlm{\nqn}{\lqn}{\mqn}.
  \end{aligned}
\end{equation}
Das Symbol $\delta_{jk}$ ist eins für $j = k$ und sonst null. Beide Aussagen sind für
die hier ausgeschriebenen Funktionen mit exakter Arithmetik nachgerechnet.

Zwei der fünf Funktionen sind komplexwertig, $\psinlm{2}{1}{1}$ und $\psinlm{2}{1}{-1}$.
Sie unterscheiden sich im Drehsinn der Phase $e^{\pm i\phi}$ und, wegen der Konvention
von Condon und Shortley, im Vorzeichen:
\begin{equation}\label{c11.eq.cs}\claim{c11.cl.cs}{numerically-checked}
  \psinlm{2}{1}{-1} = -\,\overline{\psinlm{2}{1}{1}},
  \qquad
  \abs{\psinlm{2}{1}{-1}}^2 = \abs{\psinlm{2}{1}{1}}^2 .
\end{equation}
Die Dichten der beiden Zustände sind also gleich. Ein Bild der Dichte kann sie nicht
unterscheiden (Abbildung~\ref{c11.fig.schnitt}); der Unterschied steckt in der Phase
(Abbildung~\ref{c11.fig.phase}).

\begin{tpcfigure}{c11.fig.phase}
  \demo{c11.demo.phase}
  \caption{Vorzeichen und Phase der 2p-Wellenfunktionen für $\Znuc = 1$; die Grauwerte
    zeigen wieder $\abs{\psi}^2$. Links $\psinlm{2}{1}{0}$ in der Ebene $y = 0$: Die
    Funktion ist reell, oben positiv, unten negativ. Rechts $\psinlm{2}{1}{1}$ in der
    Ebene $z = 0$: Der Betrag hängt nur vom Abstand zur Mitte ab, die Phase
    $\arg\psinlm{2}{1}{1} = \phi + \pi$ nur vom Winkel $\phi$; an den vier Strahlen
    steht ihr Wert, und der Pfeil zeigt, in welcher Richtung sie wächst. Für
    $\psinlm{2}{1}{-1}$ läuft die Phase andersherum.}
\end{tpcfigure}

\begin{remark}
  Die üblichen Orbitalbilder zeigen $\abs{\psi}^2$. Dabei geht Information verloren:
  Multipliziert man $\psi$ an jedem Punkt mit einer Zahl vom Betrag eins, die vom Ort
  abhängen darf, so ändert sich die Dichte nicht, die Wellenfunktion aber schon. Zu
  einem Bild der Dichte gehört deshalb immer die Angabe, welcher Zustand gemeint ist.
\end{remark}

Für $\nqn = 3$ und darüber geht es nach demselben Muster weiter\zoomref{c11.v3}.

\begin{vertiefung}[id=c11.v3,
    titel={Die Wellenfunktionen für $\nqn = 3$},
    wozu={Der Kern schreibt die Zustände bis $\nqn = 2$ aus. Hier stehen die neun
      Zustände zu $\nqn = 3$, darunter die ersten d-Zustände.},
    voraussetzung={Kern~\ref{c11.s5}; die allgemeine Formel aus
      Vertiefung~\ref{c11.v2} hilft, ist aber nicht nötig. Nicht nötig für die
      Abschnitte~\ref{c11.s6} und \ref{c11.s7}.},
    zurueck={Mit wachsendem $\nqn$ kommen Knoten hinzu, und die Funktionen reichen
      weiter nach außen. Die Bauart bleibt: Radialfunktion mal Funktion der Winkel.}]
  Für $\lqn = 2$ brauchen wir fünf weitere Kugelflächenfunktionen, wieder in der
  Condon-Shortley-Konvention aus \defref{D67}:
  \begin{align}
    \Ylm{2}{0}(\theta, \phi) &= \sqrt{\frac{5}{16\pi}}\,\bigl(3\cos^2\theta - 1\bigr),
      \label{c11.eq.y20}\claim{c11.cl.y20}{numerically-checked} \\
    \Ylm{2}{\pm 1}(\theta, \phi) &= \mp\sqrt{\frac{15}{8\pi}}\,\sin\theta\,\cos\theta\,e^{\pm i\phi},
      \label{c11.eq.y2pm1}\claim{c11.cl.y2pm1}{numerically-checked} \\
    \Ylm{2}{\pm 2}(\theta, \phi) &= \sqrt{\frac{15}{32\pi}}\,\sin^2\theta\,e^{\pm 2i\phi}.
      \label{c11.eq.y2pm2}\claim{c11.cl.y2pm2}{numerically-checked}
  \end{align}

  \begin{tpcfigure}{c11.fig.schnitt3}
    \caption{Die Dichte $\abs{\psi}^2$ der Relativkoordinate für die d-Zustände mit $\nqn = 3$, als Schnitt durch die Ebene $y = 0$ und für $\Znuc = 1$, in derselben Darstellung wie Abbildung~\ref{c11.fig.schnitt}; alle drei Felder benutzen eine gemeinsame lineare Grauskala. Die Achsen sind Komponenten von $\posRel = \posE - \posP$.}
  \end{tpcfigure}

  Die drei Radialfunktionen zu $\nqn = 3$ sind
  \begin{align}
    \Rnl{3}{0}(\rad) &= \frac{2}{81\sqrt{3}}\left(\frac{\Znuc}{\aRed}\right)^{3/2}
      \left(27 - 18\,\frac{\Znuc\rad}{\aRed} + 2\,\frac{\Znuc^2\rad^2}{\aRed^2}\right)
      e^{-\Znuc\rad/(3\aRed)},
      \label{c11.eq.r30}\claim{c11.cl.r30}{numerically-checked} \\
    \Rnl{3}{1}(\rad) &= \frac{4}{81\sqrt{6}}\left(\frac{\Znuc}{\aRed}\right)^{3/2}
      \frac{\Znuc\rad}{\aRed}\left(6 - \frac{\Znuc\rad}{\aRed}\right)
      e^{-\Znuc\rad/(3\aRed)},
      \label{c11.eq.r31}\claim{c11.cl.r31}{numerically-checked} \\
    \Rnl{3}{2}(\rad) &= \frac{4}{81\sqrt{30}}\left(\frac{\Znuc}{\aRed}\right)^{3/2}
      \frac{\Znuc^2\rad^2}{\aRed^2}\,e^{-\Znuc\rad/(3\aRed)},
      \label{c11.eq.r32}\claim{c11.cl.r32}{numerically-checked}
  \end{align}
  und die Wellenfunktionen $\psinlm{3}{\lqn}{\mqn} = \Rnl{3}{\lqn}\,\Ylm{\lqn}{\mqn}$
  lauten ausgeschrieben
  \begin{align}
    \psinlm{3}{0}{0}(\rad, \theta, \phi) &= \frac{1}{81\sqrt{3\pi}}
      \left(\frac{\Znuc}{\aRed}\right)^{3/2}
      \left(27 - 18\,\frac{\Znuc\rad}{\aRed} + 2\,\frac{\Znuc^2\rad^2}{\aRed^2}\right)
      e^{-\Znuc\rad/(3\aRed)},
      \label{c11.eq.psi300}\claim{c11.cl.psi300}{numerically-checked} \\
    \psinlm{3}{1}{0}(\rad, \theta, \phi) &= \frac{2}{81\sqrt{2\pi}}
      \left(\frac{\Znuc}{\aRed}\right)^{3/2}
      \frac{\Znuc\rad}{\aRed}\left(6 - \frac{\Znuc\rad}{\aRed}\right)
      e^{-\Znuc\rad/(3\aRed)}\,\cos\theta,
      \label{c11.eq.psi310}\claim{c11.cl.psi310}{numerically-checked} \\
    \psinlm{3}{1}{\pm 1}(\rad, \theta, \phi) &= \mp\frac{1}{81\sqrt{\pi}}
      \left(\frac{\Znuc}{\aRed}\right)^{3/2}
      \frac{\Znuc\rad}{\aRed}\left(6 - \frac{\Znuc\rad}{\aRed}\right)
      e^{-\Znuc\rad/(3\aRed)}\,\sin\theta\,e^{\pm i\phi},
      \label{c11.eq.psi31pm}\claim{c11.cl.psi31pm}{numerically-checked} \\
    \psinlm{3}{2}{0}(\rad, \theta, \phi) &= \frac{1}{81\sqrt{6\pi}}
      \left(\frac{\Znuc}{\aRed}\right)^{3/2}
      \frac{\Znuc^2\rad^2}{\aRed^2}\,e^{-\Znuc\rad/(3\aRed)}\,
      \bigl(3\cos^2\theta - 1\bigr),
      \label{c11.eq.psi320}\claim{c11.cl.psi320}{numerically-checked} \\
    \psinlm{3}{2}{\pm 1}(\rad, \theta, \phi) &= \mp\frac{1}{81\sqrt{\pi}}
      \left(\frac{\Znuc}{\aRed}\right)^{3/2}
      \frac{\Znuc^2\rad^2}{\aRed^2}\,e^{-\Znuc\rad/(3\aRed)}\,
      \sin\theta\,\cos\theta\,e^{\pm i\phi},
      \label{c11.eq.psi32pm1}\claim{c11.cl.psi32pm1}{numerically-checked} \\
    \psinlm{3}{2}{\pm 2}(\rad, \theta, \phi) &= \frac{1}{162\sqrt{\pi}}
      \left(\frac{\Znuc}{\aRed}\right)^{3/2}
      \frac{\Znuc^2\rad^2}{\aRed^2}\,e^{-\Znuc\rad/(3\aRed)}\,
      \sin^2\theta\,e^{\pm 2i\phi}.
      \label{c11.eq.psi32pm2}\claim{c11.cl.psi32pm2}{numerically-checked}
  \end{align}
  Die sechs Zeilen zählen neun Zustände auf. Zusammen mit den fünf Zuständen des Kerns
  sind es vierzehn; für alle vierzehn sind Normierung, Orthogonalität und die
  Eigenwertgleichung aus \eqref{c11.eq.orthonormal} exakt nachgerechnet.
  $\Rnl{3}{0}$ hat zwei Knoten, bei
  $\Znuc\rad/\aRed = (9 \mp 3\sqrt{3})/2$, den Nullstellen des quadratischen Faktors;
  $\Rnl{3}{1}$ hat einen bei $\Znuc\rad/\aRed = 6$, und $\Rnl{3}{2}$ hat keinen
  (Abbildung~\ref{c11.fig.radial3}).

  \begin{tpcfigure}{c11.fig.radial3}
    \caption{Oben die Radialfunktionen \eqref{c11.eq.r30} bis \eqref{c11.eq.r32} für
      $\Znuc = 1$, unten die zugehörigen radialen Aufenthaltsdichten
      $\Pnl{3}{\lqn}(\rad) = \rad^2\,\Rnl{3}{\lqn}(\rad)^2$ aus \defref{D72} in
      Abschnitt~\ref{c11.s7}. Durchgezogen 3s, gestrichelt 3p,
      strichpunktiert 3d. Je kleiner $\lqn$, desto mehr Knoten.}
  \end{tpcfigure}
\end{vertiefung}


\section{Spektrallinien}\label{c11.s6}

Für den spektralen Vergleich der Formel \eqref{c11.eq.energien} mit dem Labor verwenden
wir Differenzen der inneren Niveaus.

\begin{merksatz}[Spektraler Vergleich]
  Für den hier vorgenommenen spektralen Vergleich verwenden wir Differenzen der
  inneren Niveaus. Die folgende zusätzliche Photonenergie-Konvention verbindet sie
  mit Wellenlängen.
\end{merksatz}

Ein Atom kann von einem Zustand der Energie $\En{\nqn}$ in einen der niedrigeren Energie
$\En{\nqn'}$ übergehen und dabei Licht aussenden; die Energie eines Lichtquants wird mit
der Differenz $\En{\nqn} - \En{\nqn'}$ verglichen. Unser Modell enthält kein Lichtfeld.
Dass die Differenz als Licht erscheint, ist deshalb hier eine Annahme, die zum Modell
hinzukommt, und der Rückstoß, den das Atom beim Aussenden erfährt, ist dabei
vernachlässigt. Unter dieser Annahme gehört zu jedem Paar $\nqn > \nqn'$ eine
Wellenlänge $\lambda$ im Vakuum. Für Wasserstoff, $\Znuc = 1$, ist
\begin{equation}\label{c11.eq.linie}\claim{c11.cl.linie}{derived}
  \begin{aligned}
    \En{\nqn} - \En{\nqn'} &= \RydRed\,\Bigl(\frac{1}{\nqn'^2} - \frac{1}{\nqn^2}\Bigr), \\
    \lambda &= \frac{\hPlanck\,\clight}{\En{\nqn} - \En{\nqn'}}
             = \frac{1 + \massE/\massP}{\Rinf}\,
               \Bigl(\frac{1}{\nqn'^2} - \frac{1}{\nqn^2}\Bigr)^{-1} .
  \end{aligned}
\end{equation}
Hier ist $\hPlanck = 2\pi\hbar$ das Plancksche Wirkungsquantum, $\clight$ die
Lichtgeschwindigkeit und $\Rinf$ die Rydberg-Konstante
\srcref{codata-2022}{Zeile „Rydberg constant“}, mit der man die Rydberg-Energie als
$\Ryd = \hPlanck\,\clight\,\Rinf$ schreibt. Der Faktor $1 + \massE/\massP$ kommt von
der reduzierten Masse; er ist die Spur, die das Proton im Spektrum hinterlässt.

Die Übergänge, die bei $\nqn' = 1$ enden, bilden die Lyman-Serie, die bei $\nqn' = 2$
enden, die Balmer-Serie (Abbildung~\ref{c11.fig.termschema}). Die jeweils längste
Wellenlänge gehört zum kleinsten Sprung. Unser Modell liefert für den Übergang
$2 \to 1$ die Wellenlänge $\num[c11.cl.lyman]{121{,}568}{nm}$ und für den Übergang
$3 \to 2$ die Wellenlänge $\num[c11.cl.balmer]{656{,}470}{nm}$. Die zugehörigen
Energiedifferenzen stehen im Verhältnis
\begin{equation}\label{c11.eq.verhaeltnis}\claim{c11.cl.verhaeltnis}{derived}
  \frac{\En{2} - \En{1}}{\En{3} - \En{2}}
  = \frac{1 - \tfrac14}{\tfrac14 - \tfrac19}
  = \frac{3/4}{5/36} = \frac{27}{5} .
\end{equation}

\begin{tpcfigure}{c11.fig.termschema}
  \caption{Das Termschema der Relativbewegung: die Energien $\En{\nqn}$ aus
    \eqref{c11.eq.energien} für $\Znuc = 1$ in Einheiten von $\RydRed$. Die Niveaus
    häufen sich bei $E = 0$. Die Pfeile sind Übergänge; ihre Länge ist maßstäblich die
    Energiedifferenz. Durchgezogen: Lyman-Serie, $\nqn \to 1$. Gestrichelt:
    Balmer-Serie, $\nqn \to 2$.}
\end{tpcfigure}

\begin{tpcfigure}{c11.fig.spektrum}
  \caption{Dieselben Übergänge als Wellenlängen nach \eqref{c11.eq.linie}, auf einer
    logarithmischen Achse. Jeder Strich ist eine Linie; durchgezogen die Lyman-Serie
    ($\nqn = 2, \ldots, 12$), gestrichelt die Balmer-Serie ($\nqn = 3, \ldots, 12$). So
    wie sich die Niveaus bei $E = 0$ häufen, drängen sich die Linien jeder Serie an
    ihrer Seriengrenze, der Wellenlänge für $\nqn \to \infty$. Zwischen den beiden
    Serien liegt keine Linie von einer der beiden.}
\end{tpcfigure}

Wie ernst darf man diese Zahlen nehmen? Sie sind exakte Folgerungen aus dem Modell und
aus gemessenen Naturkonstanten\zoomref{c11.v4}. Das Modell selbst lässt aber weg: die
Feinstruktur, die Lamb-Verschiebung, die endliche Ausdehnung des Kerns, den Spin, äußere
Felder, die Dynamik der Abstrahlung und den Rückstoß. Die Wellenlängen sind deshalb
Aussagen über das spinlose Coulomb-Modell und keine Vorhersagen auf der Genauigkeit, mit
der man wirkliche Spektrallinien auflösen kann. Auch darüber, welche Übergänge wie
häufig vorkommen, sagt unser Hamiltonoperator nichts.

\begin{vertiefung}[id=c11.v4,
    titel={Wie genau sind die Zahlen?},
    wozu={Der Kern gibt Energien und Wellenlängen auf wenige Stellen an. Hier stehen
      die Eingangsdaten, die vollen Ergebnisse und die Auskunft, wie viele ihrer Stellen
      etwas bedeuten.},
    voraussetzung={Kern~\ref{c11.s4} und \ref{c11.s6}. Nicht nötig für
      Abschnitt~\ref{c11.s7}.},
    zurueck={Die Rechnung ist auf viele Stellen genau. Was diese Stellen physikalisch
      bedeuten, entscheidet das Modell, nicht die Arithmetik.}]
  Die übernommenen Naturkonstanten stammen aus der Liste der empfohlenen Werte von
  2022 (CODATA), aus der schon die Rydberg-Konstante in \eqref{c11.eq.linie} kam.
  Abgeleitete Modellwerte verwenden diese Eingangsdaten; illustrative Parameter sind
  gewählte Werte. Tabelle~\ref{c11.tab.codata} zeigt die benutzten Zeilen. Die Werte von $\hPlanck$, $\clight$ und $\echarge$ sind
  dort als exakt gekennzeichnet; die anderen tragen eine Standardunsicherheit.

  \begin{tpctable}{c11.tab.codata}
    \caption{Eingangsdaten: Wert und Standardunsicherheit, wie sie in der Liste stehen.}
    \begin{tabular}{lll}
      \toprule
      Größe & Wert & Standardunsicherheit \\
      \midrule
      $\massE/\massP$ & $\num[c11.cl.q]{5{,}446170214889 \cdot 10^{-4}}{}$
        & $\num[c11.cl.q]{9{,}4 \cdot 10^{-15}}{}$ \\
      $\aBohr$ & $\num[c11.cl.a0]{5{,}29177210544 \cdot 10^{-11}}{m}$
        & $\num[c11.cl.a0]{8{,}2 \cdot 10^{-21}}{m}$ \\
      $\Rinf$ & $\num[c11.cl.rinf]{10973731{,}568157}{m^{-1}}$
        & $\num[c11.cl.rinf]{1{,}2 \cdot 10^{-5}}{m^{-1}}$ \\
      $\hPlanck\,\clight\,\Rinf$ & $\num[c11.cl.ry]{13{,}605693122990}{eV}$
        & $\num[c11.cl.ry]{1{,}5 \cdot 10^{-11}}{eV}$ \\
      $\hPlanck$ & $\num[c11.cl.h]{6{,}62607015 \cdot 10^{-34}}{J\,Hz^{-1}}$ & exakt \\
      $\clight$ & $\num[c11.cl.c]{299792458}{m\,s^{-1}}$ & exakt \\
      $\echarge$ & $\num[c11.cl.e]{1{,}602176634 \cdot 10^{-19}}{C}$ & exakt \\
      \bottomrule
    \end{tabular}
  \end{tpctable}

  Aus diesen Zeilen folgen die Zahlen des Kapitels durch
  $\massRed/\massE = 1/(1 + \massE/\massP)$, durch \eqref{c11.eq.bohr} und durch
  \eqref{c11.eq.linie}. Tabelle~\ref{c11.tab.ergebnisse} zeigt sie mit allen Stellen,
  die die Eingangsdaten hergeben.

  \begin{tpctable}{c11.tab.ergebnisse}
    \caption{Ergebnisse des Modells für $\Znuc = 1$ und die Skala ihrer Unsicherheit
      nach \eqref{c11.eq.fehler}, aufgerundet.}
    \begin{tabular}{lll}
      \toprule
      Größe & Wert & Unsicherheitsskala \\
      \midrule
      $\massRed/\massE$ & $\num[c11.cl.mu-ratio]{0{,}999455679424762}{}$
        & $\num[c11.cl.mu-ratio]{9{,}4 \cdot 10^{-15}}{}$ \\
      $\aRed$ & $\num[c11.cl.amu]{5{,}29465409460 \cdot 10^{-11}}{m}$
        & $\num[c11.cl.amu]{8{,}21 \cdot 10^{-21}}{m}$ \\
      $\En{1}$ mit $\massRed$ & $\num[c11.cl.e1-mu]{-13{,}598287264283}{eV}$
        & $\num[c11.cl.e1-mu]{1{,}50 \cdot 10^{-11}}{eV}$ \\
      $\En{1}$ mit $\massE$ & $\num[c11.cl.e1-me]{-13{,}605693122991}{eV}$
        & $\num[c11.cl.e1-me]{1{,}49 \cdot 10^{-11}}{eV}$ \\
      $\lambda$ für $2 \to 1$ & $\num[c11.cl.lyman]{121{,}5684456173}{nm}$
        & $\num[c11.cl.lyman]{1{,}35 \cdot 10^{-10}}{nm}$ \\
      $\lambda$ für $3 \to 2$ & $\num[c11.cl.balmer]{656{,}4696063332}{nm}$
        & $\num[c11.cl.balmer]{7{,}25 \cdot 10^{-10}}{nm}$ \\
      \bottomrule
    \end{tabular}
  \end{tpctable}

  \begin{tpcfigure}{c11.fig.unsicherheit}
    \caption{Die Unsicherheit jeder Zeile der Tabellen~\ref{c11.tab.codata} und \ref{c11.tab.ergebnisse}, geteilt durch den Betrag des Werts, auf einer logarithmischen Achse. Durchgezogen die Eingangsdaten, gestrichelt die Ergebnisse. Am ungenauesten sind die Längen $\aBohr$ und $\aRed$; die Energien und Wellenlängen erben die Genauigkeit der Rydberg-Konstante.}
  \end{tpcfigure}

  Die Unsicherheitsskala einer Größe $f$, die aus Eingangswerten $c_1, c_2, \ldots$ mit
  den Standardunsicherheiten $u(c_1), u(c_2), \ldots$ berechnet ist, bilden wir so:
  \begin{equation}\label{c11.eq.fehler}\claim{c11.cl.fehler}{derived}
    \sum_i \Bigl\lvert \frac{\partial f}{\partial c_i} \Bigr\rvert\, u(c_i) .
  \end{equation}
  Das ist eine Schranke erster Ordnung, die gilt, gleichgültig wie die Unsicherheiten
  der Eingangswerte zusammenhängen; eine fortgepflanzte Standardunsicherheit im
  strengen Sinn ist es nicht, denn dafür müsste man diese Zusammenhänge kennen.

  Die Tabellen sagen zweierlei. Die Rechnung selbst ist kein Engpass: Die Eingangsdaten
  begrenzen die Ergebnisse erst weit hinter dem Komma. Und gerade deshalb darf man die
  vielen Stellen nicht für Physik halten. Die Wellenlängen sind die des Modells aus
  Abschnitt~\ref{c11.s1}; was es weglässt, steht am Ende von Abschnitt~\ref{c11.s6}.
\end{vertiefung}

\section{Der Abstand von Elektron und Proton}\label{c11.s7}

Die Wellenfunktionen aus Abschnitt~\ref{c11.s5} sagen, wie der Abstandsvektor $\posRel$
verteilt ist. Oft will man weniger wissen: nur, wie groß der Abstand
$\rad = \abs{\posRel}$ zwischen Elektron und Proton ist, gleichgültig in welcher
Richtung.

Für \defref{D72} wählen wir die hier konstruierten Radialfunktionen reell; ihre
Quadrate sind deshalb zugleich ihre Betragsquadrate.

\tpcdef{D72}

Der Name ist berechtigt: Integriert man $\abs{\psinlm{\nqn}{\lqn}{\mqn}}^2$ über eine
Kugelschale, so liefern die Winkel den Faktor eins, weil die Kugelflächenfunktionen
normiert sind, und es bleibt
\begin{equation}\label{c11.eq.schale}\claim{c11.cl.schale}{derived}
  \int_{r_1 \leq \abs{\posRel} \leq r_2} \abs{\psinlm{\nqn}{\lqn}{\mqn}(\posRel)}^2\,\dd^3\posRel
  = \int_{r_1}^{r_2} \Pnl{\nqn}{\lqn}(\rad)\,\dd\rad .
\end{equation}
Also ist $\Pnl{\nqn}{\lqn}(\rad)\,\dd\rad$ die Wahrscheinlichkeit, den Abstand zwischen
$\rad$ und $\rad + \dd\rad$ zu finden, und zwar für jedes $\mqn$ dieselbe. Mit den
Radialfunktionen \eqref{c11.eq.r10} bis \eqref{c11.eq.r21} und $x = \Znuc\rad/\aRed$
ergibt sich
\begin{equation}\label{c11.eq.pnl}\claim{c11.cl.pnl}{derived}
  \Pnl{1}{0} = \frac{\Znuc}{\aRed}\,4x^2\,e^{-2x},
  \qquad
  \Pnl{2}{0} = \frac{\Znuc}{\aRed}\,\frac{x^2(2 - x)^2}{8}\,e^{-x},
  \qquad
  \Pnl{2}{1} = \frac{\Znuc}{\aRed}\,\frac{x^4}{24}\,e^{-x} .
\end{equation}

\begin{tpcfigure}{c11.fig.dichte}
  \demo{c11.demo.radial}
  \caption{Die radialen Aufenthaltsdichten \eqref{c11.eq.pnl} für $\Znuc = 1$.
    Durchgezogen 1s, gestrichelt 2s, strichpunktiert 2p. Die gepunktete Linie markiert
    das Maximum von $\Pnl{1}{0}$ bei $\rad = \aRed$, der Punkt den Knoten von
    $\Pnl{2}{0}$ bei $\rad = 2\aRed$. Die kurzen Striche unter der Achse sind die
    mittleren Abstände \eqref{c11.eq.mittel}. Alle Kurven beginnen bei null, auch die
    der s-Zustände: Der Faktor $\rad^2$ kommt vom Volumen der Kugelschale.}
\end{tpcfigure}

Drei Dinge liest man an Abbildung~\ref{c11.fig.dichte} ab.

Erstens der \emph{wahrscheinlichste Abstand} im Zustand 1s. Die Ableitung von
$x^2 e^{-2x}$ ist $2x(1 - x)\,e^{-2x}$; sie verschwindet für $x > 0$ nur bei $x = 1$.
Also hat $\Pnl{1}{0}$ sein Maximum bei
\begin{equation}\label{c11.eq.maximum}\claim{c11.cl.maximum}{numerically-checked}
  \rad = \frac{\aRed}{\Znuc} .
\end{equation}
Das gibt dem Bohrschen Radius eine anschauliche Bedeutung: Für Wasserstoff ist $\aRed$
der wahrscheinlichste Abstand zwischen Elektron und Proton im Zustand 1s. Die
Wellenfunktion $\psinlm{1}{0}{0}$ selbst ist dagegen bei $\rad = 0$ am größten. Das ist
kein Widerspruch: $\abs{\psi}^2$ ist eine Dichte pro Volumen, $\Pnl{1}{0}$ eine Dichte
pro Abstand, und zu einem größeren Abstand gehört eine größere Kugelschale.

\begin{tpcfigure}{c11.fig.kugelschale}
  \caption{Warum die radiale Aufenthaltsdichte von 1s bei null beginnt, für $\Znuc = 1$. Gestrichelt: die skalierte Dichte $\pi\aRed^3\abs{\psinlm{1}{0}{0}}^2 = e^{-2x}$ pro Volumen, am Ursprung am größten. Dünn: $\rad^2/\aRed^2 = x^2$, die Kugeloberfläche geteilt durch $4\pi\aRed^2$. Dick: $\aRed\Pnl{1}{0} = 4x^2 e^{-2x}$, die skalierte Dichte pro Abstand, also das Vierfache des Produkts der dünnen und der gestrichelten Kurve; ihr Maximum liegt bei $\rad = \aRed$. Hier ist $x = \rad/\aRed$.}
\end{tpcfigure}

Zweitens die \emph{Knoten}. $\Pnl{2}{0}$ verschwindet bei $x = 2$, also bei
$\rad = 2\aRed/\Znuc$, weil dort $\Rnl{2}{0}$ das Vorzeichen wechselt. Im Zustand 2s ist
die Dichte auf der ganzen Kugel mit diesem Radius null. Die Zustände 1s und 2p haben
keinen solchen Knoten.

Drittens der \emph{mittlere Abstand} $\expval{\rad} = \int_0^\infty \rad\,
\Pnl{\nqn}{\lqn}(\rad)\,\dd\rad$. Für den Zustand 1s ist das
$(\aRed/\Znuc) \cdot 4\int_0^\infty x^3 e^{-2x}\,\dd x = (\aRed/\Znuc) \cdot 4 \cdot
3!/2^4$. Dieselbe Rechnung für die beiden anderen Dichten ergibt zusammen
\begin{equation}\label{c11.eq.mittel}\claim{c11.cl.mittel}{derived}
  \expval{\rad}_{10} = \frac{3}{2}\,\frac{\aRed}{\Znuc},
  \qquad
  \expval{\rad}_{20} = 6\,\frac{\aRed}{\Znuc},
  \qquad
  \expval{\rad}_{21} = 5\,\frac{\aRed}{\Znuc} .
\end{equation}
Die Indizes sind $\nqn$ und $\lqn$. Der mittlere Abstand ist größer als der
wahrscheinlichste, weil die Dichte nach außen langsam ausläuft. Und obwohl 2s und 2p
dieselbe Energie haben, ist der mittlere Abstand verschieden. Für alle $\nqn$ und
$\lqn$ gibt es dafür eine geschlossene Formel\zoomref{c11.v5}.

Damit sind die konstruierten gebundenen Zustände und ihre Dichten beschrieben. Die
Bewegung des Schwerpunkts ist frei, die konstruierte Familie hat die Energien
\eqref{c11.eq.energien} und die Wellenfunktionen aus Abschnitt~\ref{c11.s5}, und was man
im Labor von einem Atom sieht, verlangt beide Hälften (Abbildung~\ref{c11.fig.ueberblick}).

\begin{tpcfigure}{c11.fig.ueberblick}
  \caption{Der Weg des Kapitels als Schema. Aus dem Zweiteilchenproblem werden zwei unabhängige Probleme; die Relativbewegung liefert Energien und Wellenfunktionen. Mit dem Labor verglichen werden die Differenzen der Energien, als Spektrallinien, und Dichten, in die der Zustand des Schwerpunkts eingeht.}
\end{tpcfigure}


\begin{vertiefung}[id=c11.v5,
    titel={Momente und Verteilungsfunktion},
    wozu={Der Kern berechnet mittlere Abstände für $\nqn \leq 2$ einzeln. Hier steht
      eine Formel für alle Zustände, dazu die Wahrscheinlichkeit, den Abstand unterhalb
      eines gegebenen Radius zu finden.},
    voraussetzung={Kern~\ref{c11.s7}; für den allgemeinen Fall die Radialfunktionen aus
      Vertiefung~\ref{c11.v2}.},
    zurueck={Mittlere Abstände und eingeschlossene Wahrscheinlichkeiten gibt es für
      jeden Zustand als Formel, ohne numerische Integration.}]
  Wir messen Längen wieder mit $x = \Znuc\rad/\aRed$ und schreiben $u(x)$ für die auf
  $\int_0^\infty u^2\,\dd x = 1$ normierte reelle Lösung von \eqref{c11.eq.dimlos} zu
  $\varepsilon = -1/\nqn^2$. Die Momente sind $I_s = \int_0^\infty x^s\,u(x)^2\,\dd x$.
  Multipliziert man \eqref{c11.eq.dimlos} einmal mit $x^s u$ und einmal mit
  $x^{s + 1} u'$ und integriert partiell, so erhält man zwei Gleichungen, aus denen man
  das Integral über $x^s (u')^2$ eliminieren kann. Für $s = 0, 1, 2$ verschwinden die
  Randterme, und es bleibt
  \begin{equation}\label{c11.eq.kramers}\claim{c11.cl.kramers}{derived}
    s\,\Bigl(\frac{s^2 - 1}{4} - \lqn(\lqn + 1)\Bigr)\,I_{s - 2}
    + (2s + 1)\,I_{s - 1} - \frac{s + 1}{\nqn^2}\,I_s = 0 .
  \end{equation}
  Mit $I_0 = 1$ liefern $s = 0$, $s = 1$ und $s = 2$ der Reihe nach
  \begin{equation}\label{c11.eq.momente}\claim{c11.cl.momente}{derived}
    \begin{aligned}
      \Bigl\langle \frac{1}{\rad} \Bigr\rangle &= \frac{\Znuc}{\aRed\,\nqn^2},
      \qquad
      \expval{\rad} = \frac{\aRed}{2\Znuc}\,\bigl(3\nqn^2 - \lqn(\lqn + 1)\bigr), \\
      \expval{\rad^2} &= \frac{\aRed^2\,\nqn^2}{2\Znuc^2}\,
        \bigl(5\nqn^2 + 1 - 3\lqn(\lqn + 1)\bigr).
    \end{aligned}
  \end{equation}
  Für $\nqn \leq 2$ sind das die Werte aus \eqref{c11.eq.mittel}. Bei festem $\nqn$ ist
  der mittlere Abstand umso kleiner, je größer $\lqn$ ist.

  Für die zweite Frage definieren wir die Verteilungsfunktion.

  \tpcdef{D73}

  Weil $\Pnl{\nqn}{\lqn}$ ein Polynom mal $e^{-2x/\nqn}$ ist, lässt sich das Integral
  durch wiederholtes partielles Integrieren geschlossen ausführen:
  \begin{equation}\label{c11.eq.cdf}\claim{c11.cl.cdf}{derived}
    \Fnl{\nqn}{\lqn}(\rad) = 1 - e^{-2x/\nqn}\,\Qnl{\nqn}{\lqn}(x),
    \qquad
    x = \frac{\Znuc\rad}{\aRed},
  \end{equation}
  mit einem Polynom $\Qnl{\nqn}{\lqn}$. Für die Zustände bis $\nqn = 3$ ist
  \begin{equation}\label{c11.eq.qpolynome}\claim{c11.cl.qpolynome}{derived}
    \begin{aligned}
      \Qnl{1}{0} &= 1 + 2x + 2x^2, \\
      \Qnl{2}{0} &= 1 + x + \frac{x^2}{2} + \frac{x^4}{8}, \\
      \Qnl{2}{1} &= 1 + x + \frac{x^2}{2} + \frac{x^3}{6} + \frac{x^4}{24}, \\
      \Qnl{3}{0} &= 1 + \frac{2x}{3} + \frac{2x^2}{9} + \frac{4x^4}{81}
                    - \frac{8x^5}{729} + \frac{8x^6}{6561}, \\
      \Qnl{3}{1} &= 1 + \frac{2x}{3} + \frac{2x^2}{9} + \frac{4x^3}{81}
                    + \frac{2x^4}{243} - \frac{4x^5}{2187} + \frac{4x^6}{6561}, \\
      \Qnl{3}{2} &= 1 + \frac{2x}{3} + \frac{2x^2}{9} + \frac{4x^3}{81}
                    + \frac{2x^4}{243} + \frac{4x^5}{3645} + \frac{4x^6}{32805}.
    \end{aligned}
  \end{equation}

  \begin{tpcfigure}{c11.fig.verteilung}
    \caption{Die Verteilungsfunktionen $\Fnl{1}{0}$, $\Fnl{2}{0}$ und $\Fnl{2}{1}$ nach
      \eqref{c11.eq.cdf} für $\Znuc = 1$: die Wahrscheinlichkeit, den Abstand kleiner
      als $\rad$ zu finden. Am Knoten von 2s, dem markierten Punkt bei
      $\rad = 2\aRed$, hat $\Fnl{2}{0}$ eine waagerechte Tangente.}
  \end{tpcfigure}

  Mit $\Fnl{\nqn}{\lqn}$ lässt sich sagen, was die Oberfläche in einem Orbitalbild
  bedeuten soll, und dafür gibt es zwei verschiedene Verabredungen. Die Kugel mit dem
  Radius $r_p$, für den $\Fnl{\nqn}{\lqn}(r_p) = p$ ist, schließt die Wahrscheinlichkeit
  $p$ ein, für jedes $\mqn$; sie ist aber im Allgemeinen keine Fläche gleicher Dichte.
  Eine Fläche gleicher Dichte, die die Wahrscheinlichkeit $p$ einschließt, erhält man
  aus einer Schwelle $c_p$ mit
  \begin{equation}\label{c11.eq.isoflaeche}\claim{c11.cl.isoflaeche}{derived}
    \int_{\abs{\psi}^2 \geq c_p} \abs{\psi(\posRel)}^2\,\dd^3\posRel = p ;
  \end{equation}
  das eingeschlossene Gebiet kann aus mehreren Teilen bestehen. Für den Zustand 1s, in
  dem die Dichte mit dem Abstand streng abnimmt, fallen beide Verabredungen zusammen.
  Eine Dichteschwelle, die man nicht integriert hat, darf man nicht mit einer
  Wahrscheinlichkeit beschriften.
\end{vertiefung}

\begin{exercise}[Der Zustand 2p durch Einsetzen]\label{c11.exc.zweip}
  Zeige, dass $u(x) = x^2\,e^{-x/2}$ die Gleichung \eqref{c11.eq.dimlos} für $\lqn = 1$
  löst, und lies $\varepsilon$ ab. Vergleiche mit \eqref{c11.eq.r21}.
\end{exercise}
\begin{solution}
  Es ist $u' = (2x - \tfrac12 x^2)\,e^{-x/2}$ und
  $u'' = (2 - 2x + \tfrac14 x^2)\,e^{-x/2}$. Für $\lqn = 1$ ist
  $\lqn(\lqn + 1)/x^2 - 2/x = 2/x^2 - 2/x$, und dieser Faktor mal $u$ ist
  $(2 - 2x)\,e^{-x/2}$. Zusammen:
  $-u'' + (2/x^2 - 2/x)\,u = -\tfrac14 x^2\,e^{-x/2} = -\tfrac14 u$. Also ist
  $\varepsilon = -\tfrac14$, das ist $\En{2}$ in Einheiten von $\Znuc^2\,\RydRed$. Die
  Radialfunktion $u/x$ ist bis auf den Vorfaktor $\Rnl{2}{1}$.
\end{solution}

\begin{exercise}[Mittlerer Abstand im Zustand 2p]\label{c11.exc.mittel}
  Berechne $\expval{\rad}$ für den Zustand 2p aus $\Pnl{2}{1}$ in \eqref{c11.eq.pnl}.
\end{exercise}
\begin{solution}
  Mit $\rad = (\aRed/\Znuc)\,x$ und $\dd\rad = (\aRed/\Znuc)\,\dd x$ ist
  \begin{equation*}
    \expval{\rad} = \frac{\aRed}{\Znuc} \cdot \frac{1}{24}\int_0^\infty x^5 e^{-x}\,\dd x
    = \frac{\aRed}{\Znuc} \cdot \frac{5!}{24} = 5\,\frac{\aRed}{\Znuc},
  \end{equation*}
  wie in \eqref{c11.eq.mittel}.
\end{solution}

\begin{exercise}[Das unendlich schwere Proton]\label{c11.exc.unendlich}
  Um welchen Faktor ändern sich der Bohrsche Radius und die Energien $\En{\nqn}$, wenn
  man die reduzierte Masse $\massRed$ durch die Elektronenmasse $\massE$ ersetzt?
  \begin{enumerate}
    \item Benutze \eqref{c11.eq.bohr}.
    \item Was bedeutet das für die Wellenlängen in \eqref{c11.eq.linie}?
  \end{enumerate}
\end{exercise}
\begin{solution}
  Aus $\aRed$ wird $\aBohr = \aRed\,\massRed/\massE$. Der Radius wird also um den
  Faktor $\massRed/\massE = 1/(1 + \massE/\massP)$ kleiner. Aus $\RydRed$ wird
  $\Ryd = \RydRed\,\massE/\massRed$. Alle Energien werden also dem Betrag nach um den
  Faktor $1 + \massE/\massP$ größer.
  \begin{itemize}
    \item Die Energiedifferenzen wachsen um denselben Faktor.
    \item Die Wellenlängen werden um diesen Faktor kürzer: In \eqref{c11.eq.linie}
      fällt der Faktor $1 + \massE/\massP$ weg.
  \end{itemize}
\end{solution}


\begin{fundament}[id=c11.f1,
    titel={Selbstadjungiertheit der Hamiltonoperatoren des Wasserstoffatoms},
    wozu={Im Kern sind $\op{\HamTot}$ und $\op{\HamRel}$ Rechenvorschriften für differenzierbare
      Funktionen. Hier steht, auf welchen Definitionsbereichen sie selbstadjungierte
      Operatoren festlegen, und auf welchen Satz sich das stützt.},
    voraussetzung={Kern~\ref{c11.s1} und \ref{c11.s2}; die Begriffe Definitionsbereich
      und (wesentlich) selbstadjungiert aus \defref{D16} und \defref{D19}.},
    zurueck={Die Rechenvorschriften aus \defref{D60} und
      Theorem~\ref{c11.thm.separation} legen, auf glatten Funktionen mit kompaktem
      Träger gestartet, jeweils einen selbstadjungierten Operator fest: ihren
      Abschluss.}]
  Ein Operator ist nach \defref{D16} erst angegeben, wenn man seinen Definitionsbereich
  nennt. Wir starten mit einem kleinen Bereich, auf dem an den Rechenvorschriften nichts
  zu deuten ist: den glatten Funktionen mit kompaktem Träger, $\Ccinf{\mathbb{R}^3}$ für
  die Relativbewegung und $\Ccinf{\mathbb{R}^6}$ für das ganze Atom. Die Frage ist, ob
  die Operatoren dort wesentlich selbstadjungiert sind (\defref{D19}); dann, und nur
  davon handelt dieser Block, ist ihr Abschluss selbstadjungiert.

  \begin{tpcfigure}{c11.fig.bereiche}
    \caption{Die drei Funktionenmengen dieses Blocks als Schema: die glatten Funktionen mit kompaktem Träger, der Definitionsbereich $\dom{-\Lap}$ des Laplace-Operators und der ganze Raum $\Ltwo{\mathbb{R}^3}$. Eingetragen ist, was Korollar~\ref{c11.cor.relativ} über $\op{\HamRel}$ sagt.}
  \end{tpcfigure}

  Die Antwort hängt an einer Eigenschaft des Potentials. Man schreibt
  $\Ltwo{\mathbb{R}^3} + L^{\infty}(\mathbb{R}^3)$ für die Menge der messbaren Funktionen
  $f$ auf $\mathbb{R}^3$, die sich als Summe $f = f_1 + f_2$ einer quadratintegrierbaren
  Funktion $f_1$ und einer beschränkten Funktion $f_2$ schreiben lassen. Für solche
  Potentiale gelten zwei Sätze von Kato, die wir ohne Beweis übernehmen. Im ersten ist
  $\Lap$ der Laplace-Operator auf $\mathbb{R}^3$ und $\dom{-\Lap}$ sein
  Definitionsbereich.

  \begin{theorem}[Kato; Reed und Simon, Theorem X.15]\label{c11.thm.kato15}\claim[proof-omitted]{c11.cl.kato15}{source-quoted}
    Sei $\Vpot \in \Ltwo{\mathbb{R}^3} + L^{\infty}(\mathbb{R}^3)$ reellwertig. Dann ist
    $-\Lap + \Vpot$ wesentlich selbstadjungiert auf $\Ccinf{\mathbb{R}^3}$ und
    selbstadjungiert auf $\dom{-\Lap}$.
  \end{theorem}

  \begin{theorem}[Satz von Kato; Reed und Simon, Theorem X.16]\label{c11.thm.kato16}\claim[proof-omitted]{c11.cl.kato16}{source-quoted}
    Seien $\Vpot_1, \ldots, \Vpot_m$ reellwertige messbare Funktionen, jede in
    $\Ltwo{\mathbb{R}^3} + L^{\infty}(\mathbb{R}^3)$. Sei $\Vpot_k(y_k)$ der
    Multiplikationsoperator auf $\Ltwo{\mathbb{R}^{3n}}$, der entsteht, wenn man für
    $y_k$ drei Koordinaten des $\mathbb{R}^{3n}$ wählt. Dann ist
    $-\Lap + \sum_{k = 1}^{m} \Vpot_k(y_k)$ wesentlich selbstadjungiert auf
    $\Ccinf{\mathbb{R}^{3n}}$; dabei ist $\Lap$ der Laplace-Operator auf
    $\mathbb{R}^{3n}$.
  \end{theorem}

  Beide Sätze stehen bei \srcref{reed-simon-2}{§\,X.2, S. 165--166}; dort fehlt der Faktor
  $\hbar^2$, weil mit $\hbar = 1$ gerechnet wird. Dass das Coulombpotential die
  Voraussetzung erfüllt, rechnen wir nach.

  \begin{lemma}[Das Coulombpotential ist zulässig]\label{c11.lem.coulomb}\claim[proof-included]{c11.cl.coulomb}{derived}
    Sei $c > 0$ und $\Vpot(\posRel) = -c/\abs{\posRel}$. Für jedes $b > 0$ zerlegen wir
    $\Vpot = \Vpot_1 + \Vpot_2$, wobei $\Vpot_1$ für $\abs{\posRel} < b$ gleich $\Vpot$
    und sonst null ist. Dann gilt
    \begin{equation}\label{c11.eq.coulombsplit}
      \int \abs{\Vpot_1(\posRel)}^2\,\dd^3\posRel = 4\pi c^2 b,
      \qquad
      \abs{\Vpot_2(\posRel)} \leq \frac{c}{b}\quad\text{\tpcifenglish{for all}{für alle}}\ \posRel .
    \end{equation}
    Insbesondere liegt $\Vpot$ in $\Ltwo{\mathbb{R}^3} + L^{\infty}(\mathbb{R}^3)$.
  \end{lemma}
  \begin{proof}
    In Kugelkoordinaten ist $\dd^3\posRel = 4\pi\rad^2\,\dd\rad$ für Funktionen, die nur
    von $\rad$ abhängen. Also ist das Integral gleich
    $4\pi\int_0^b (c^2/\rad^2)\,\rad^2\,\dd\rad = 4\pi c^2 b$: Der Faktor $\rad^2$ des
    Volumenelements hebt die Singularität auf. Außerhalb der Kugel ist
    $c/\abs{\posRel} \leq c/b$ (Abbildung~\ref{c11.fig.coulombsplit}).
  \end{proof}

  \begin{tpcfigure}{c11.fig.coulombsplit}
    \caption{Die Zerlegung aus Lemma~\ref{c11.lem.coulomb}. Innerhalb des Radius $b$
      (durchgezogen) wächst $\abs{\Vpot}$ über alle Grenzen, ist aber
      quadratintegrierbar; außerhalb (gestrichelt) ist $\abs{\Vpot}$ durch $c/b$
      beschränkt.}
  \end{tpcfigure}

  \begin{corollary}[Relativbewegung]\label{c11.cor.relativ}\claim[proof-included]{c11.cl.sa-relativ}{source-backed}
    Der Operator $\op{\HamRel} = -\frac{\hbar^2}{2\massRed}\Lap + \Vpot$ mit dem
    Coulombpotential aus \defref{D60} ist wesentlich selbstadjungiert auf
    $\Ccinf{\mathbb{R}^3}$ und selbstadjungiert auf $\dom{-\Lap}$.
  \end{corollary}
  \begin{proof}
    Es ist $\op{\HamRel} = \frac{\hbar^2}{2\massRed}\,(-\Lap + W)$ mit
    $W = \frac{2\massRed}{\hbar^2}\,\Vpot$, und $W(\posRel) = -c/\abs{\posRel}$ mit
    $c = 2\Znuc/\aRed > 0$. Nach Lemma~\ref{c11.lem.coulomb} erfüllt $W$ die
    Voraussetzung von Theorem~\ref{c11.thm.kato15}. Ein reeller Faktor $c \neq 0$
    erhält die Definitionsbereiche und die wesentliche Selbstadjungiertheit:
    $(cA)^* = cA^*$ und $\overline{cA} = c\overline{A}$. Die erste Identität folgt aus
    der Definitionsgleichung des Adjungierten, die zweite aus der umkehrbaren
    Graphenabbildung $(f,g) \mapsto (f,cg)$. Der positive Faktor überträgt also beide
    Aussagen von $-\Lap + W$ auf $\op{\HamRel}$.
  \end{proof}

  Für das Coulombpotential steht die erste Aussage des Korollars bei Reed und Simon als
  Beispiel~1 direkt hinter dem Beweis von Theorem X.15.

  \begin{proposition}[Das ganze Atom]\label{c11.prop.zweikoerper}\claim[proof-included]{c11.cl.sa-zweikoerper}{source-backed}
    Der Zweikörper-Hamiltonoperator $\op{\HamTot}$ aus \defref{D60} mit dem Coulombpotential
    ist wesentlich selbstadjungiert auf $\Ccinf{\mathbb{R}^6}$.
  \end{proposition}
  \begin{proof}
    Theorem~\ref{c11.thm.kato16} verlangt den Laplace-Operator des $\mathbb{R}^6$ ohne
    Massenfaktoren und ein Potential, das von drei der Koordinaten abhängt. Beides
    stellen wir in zwei Schritten her.
    \begin{description}
      \item[Schritt 1, Massen wegskalieren] Wir führen neue Variablen ein,
        \begin{equation*}
          \boldsymbol{x}_{\mathrm{e}} = \frac{\sqrt{2\massE}}{\hbar}\,\posE,
          \qquad
          \boldsymbol{x}_{\mathrm{p}} = \frac{\sqrt{2\massP}}{\hbar}\,\posP .
        \end{equation*}
        In ihnen wird die kinetische Energie zu $-\Lap$ auf $\mathbb{R}^6$. Das Umskalieren ist, mit einem konstanten
        Normierungsfaktor versehen, eine normerhaltende, umkehrbare Abbildung, die
        $\Ccinf{\mathbb{R}^6}$ auf sich abbildet.
      \item[Schritt 2, die Relativrichtung] Im Potential steht
        $\posE - \posP = \hbar\gamma\,\boldsymbol{z}$ mit
        \begin{equation}\label{c11.eq.drehung}
          \boldsymbol{z} = \frac{1}{\gamma}\Bigl(
            \frac{\boldsymbol{x}_{\mathrm{e}}}{\sqrt{2\massE}}
            - \frac{\boldsymbol{x}_{\mathrm{p}}}{\sqrt{2\massP}}\Bigr),
          \qquad
          \gamma = \sqrt{\frac{1}{2\massE} + \frac{1}{2\massP}}
                 = \frac{1}{\sqrt{2\massRed}} .
        \end{equation}
        Die Quadrate der beiden Vorfaktoren sind $1/(2\massE\gamma^2)$ und
        $1/(2\massP\gamma^2)$. Ihre Summe ist eins. Die drei Komponenten von
        $\boldsymbol{z}$ sind deshalb Koordinaten des $\mathbb{R}^6$ längs dreier
        Einheitsrichtungen, die aufeinander senkrecht stehen. Nach einer Drehung des
        $\mathbb{R}^6$ sind sie drei der sechs Koordinaten. In ihnen ist das Potential
        $-c/\abs{\boldsymbol{z}}$ mit einem $c > 0$. Nach Lemma~\ref{c11.lem.coulomb}
        ist es also zulässig.
    \end{description}
    Theorem~\ref{c11.thm.kato16} mit $n = 2$ und $m = 1$ ist für drei der Koordinaten
    selbst formuliert; sein Beweis bei Reed und Simon beginnt mit genau dieser Drehung
    der Variablen und deckt den gedrehten Fall ab. Drehung und Umskalierung bilden
    $\Ccinf{\mathbb{R}^6}$ auf sich ab und erhalten die wesentliche Selbstadjungiertheit.
  \end{proof}

  \begin{tpcfigure}{c11.fig.drehung}
    \caption{Schritt 2 des Beweises für eine Raumrichtung, gezeichnet für $\massP : \massE = 3 : 1$. In der Ebene der skalierten Koordinaten $x_{\mathrm{e}}$, $x_{\mathrm{p}}$ ist $z$ die Koordinate längs des Einheitsvektors mit den Komponenten $\sqrt{\massP/\massTot}$ und $-\sqrt{\massE/\massTot}$. Dieser und die dazu senkrechte Richtung bilden ein gedrehtes Achsenkreuz; der Kreis ist der Einheitskreis.}
  \end{tpcfigure}

  Reed und Simon ziehen denselben Schluss für jedes Potential der Klasse
  $\Ltwo{\mathbb{R}^3} + L^{\infty}(\mathbb{R}^3)$ in einem Satz
  \srcref{reed-simon-3}{§\,XI.4, S. 55}. Beim Vergleich mit dieser Seite ist zu beachten:
  Die beiden Massen und die reduzierte Masse tragen dort andere Buchstaben (die
  reduzierte Masse heißt $m$); es ist $\hbar = 1$; die glatten Funktionen mit kompaktem
  Träger heißen $C_0^{\infty}$; und die dort $U$ genannte Abbildung macht aus einer
  Funktion der neuen Koordinaten eine Funktion der alten, sie ist also die Umkehrung
  unserer Abbildung $\op{\Uni}$ aus \eqref{c11.eq.tilde}. Die Rechnung, die zu
  Theorem~\ref{c11.thm.separation} führt, ist dort eine Übungsaufgabe.
\end{fundament}
